%% =====================================================================
%%  Capítulo 10 -- Más allá de la estabilidad asintótica: UUB, conversos,
%%  CLF, Barbalat e ISS
%%  Formato nuevo (noviembre de 2026). Integra: AdvancedStab_Notas2026.tex
%%  (notas de octubre de 2026, que ya corregían Handout7 y las diapositivas),
%%  las diapositivas AdvancedStab_2026_v4 (clase del 28 de octubre, "CLF e
%%  ISS, con conversos y Barbalat como apoyo"), las ecuaciones de LaTeXiT y
%%  los scripts MATLAB_AdvancedStab/AS01--AS15, reescritos en
%%  codigo/cap10_avanzada. Las funciones de comparación, la estabilidad
%%  uniforme y LaSalle--Yoshizawa son del capítulo 9 (por migrar): aquí se
%%  recuerda lo indispensable. Correcciones: ver \S\ref{sec:av-cierre}.
%%  Números verificados con herramientas/cap10_numeros.py y
%%  codigo/cap10_avanzada/verificarCap10.m.
%% =====================================================================

\capitulo[cap:avanzada]{Más allá de la estabilidad asintótica: UUB, conversos, CLF, Barbalat e ISS}
  {Corren tiempos raros. El arte siempre está fuera de lugar. Siempre es innecesario cuando surge. Y después, al mirar atrás, resulta que era lo único importante.}
  {Daniel Kehlmann}

\begin{objetivos}[partes/10-estabilidad-avanzada.html]
  \item Garantizar acotamiento y acotamiento último uniforme (UB, UUB) cuando no hay equilibrio, con el argumento de los dos conjuntos de nivel, y calcular la cota última que de verdad garantiza el teorema.
  \item Explicar qué entrega un teorema converso de Lyapunov según la hipótesis (exponencial, asintótica, solo estable) y usar sus cotas para probar robustez frente a perturbaciones.
  \item Decidir si una función es de control de Lyapunov, diseñar con la fórmula universal de Sontag y reconocer cuándo no existe una realimentación continua (Artstein, control pequeño, Brockett).
  \item Usar el lema de Barbalat en sistemas variantes en el tiempo, verificando la continuidad uniforme con una derivada acotada, y saber qué concluye y qué no.
  \item Probar estabilidad entrada--estado con funciones ISS-Lyapunov, robustecer un diseño con amortiguamiento $\Lie gV$ y analizar cascadas y retroalimentaciones con la condición de pequeña ganancia.
\end{objetivos}

%% ---------------------------------------------------------------------
\section{Cinco preguntas que la estabilidad asintótica no responde}
\label{sec:av-porque}
%% ---------------------------------------------------------------------

El \cref{cap:lyapunov} y el capítulo~9 respondieron una pregunta: si $\dot x=f(t,x)$ tiene un equilibrio en el origen, ¿las soluciones se quedan cerca y llegan a él? En ingeniería esa pregunta casi nunca es la única. Este capítulo responde otras cinco:
\begin{enumerate}
  \item \textbf{Sin equilibrio} (una perturbación que no se apaga), ¿qué se puede garantizar? $\to$ acotamiento último uniforme (UUB), \S\ref{sec:av-uub}.
  \item Si no encuentro una función de Lyapunov, ¿es porque no existe? $\to$ teoremas conversos, \S\ref{sec:av-conversos}.
  \item \textbf{Diseño:} ¿cómo escojo la entrada $u$ usando una $V$? $\to$ funciones de control de Lyapunov y fórmula de Sontag, \S\ref{sec:av-clf}.
  \item Si $\Vdot\le0$ en un sistema \textbf{variante en el tiempo}, ¿qué converge? $\to$ lema de Barbalat, \S\ref{sec:av-barbalat}.
  \item Con entradas: ¿entrada acotada da estado acotado? ¿Y al interconectar dos sistemas? $\to$ estabilidad entrada--estado (ISS) y pequeña ganancia, \S\ref{sec:av-iss}--\ref{sec:av-inter}.
\end{enumerate}

Hay un hilo que une las respuestas, y conviene decirlo desde el principio: \alerta{casi todas son variaciones del mismo argumento de conjuntos de nivel}. Si $V$ decrece fuera de un conjunto, las trayectorias quedan atrapadas en el menor conjunto de nivel de $V$ que contiene a ese conjunto. Con esa sola idea se prueban el teorema de UUB, el de funciones ISS-Lyapunov y, con un paso más, el de cascadas (\cref{fig:av-mapa}).

\begin{figure}[ht]
\centering
\includegraphics[width=0.95\linewidth]{cap10/t_mapa_avanzada}
\caption{Mapa del capítulo. El argumento de conjuntos de nivel da UUB, ISS y cascadas; los conversos garantizan que las funciones que se suponen existen; las CLF pasan del análisis al diseño; Barbalat cubre lo que LaSalle no puede en sistemas variantes en el tiempo.}
\label{fig:av-mapa}
\end{figure}

\notamargen{\textbf{Ruta de las clases.} Semana~11 (21 de octubre): UUB, \S\ref{sec:av-uub}. Semana~12 (26 y 28 de octubre): conversos de apoyo, CLF y Sontag, Barbalat, ISS e interconexiones, \S\ref{sec:av-conversos}--\ref{sec:av-inter}. El 28 de octubre la clase se concentró en CLF e ISS; las demás secciones se leen por cuenta propia.}

\begin{ideaclave}
\textbf{``Existe $u$'' contra ``para todo $w$''.} La letra de la entrada cambia de papel en el capítulo. En las CLF (\S\ref{sec:av-clf}), $u$ es un \emph{control que se diseña}, y la pregunta es si \emph{existe} un $u$ que haga bajar a $V$. En ISS (\S\ref{sec:av-iss}), $w$ es una \emph{entrada exógena} que no se controla (perturbación, error de actuador), y la pregunta es si el estado queda acotado \emph{para toda} $w$ acotada. En las interconexiones (\S\ref{sec:av-inter}), la entrada de un subsistema es el estado del otro. Se escribe $u$ para lo que se diseña y $w$ para lo que se sufre, como en las convenciones del curso; \citet{Kha02} escribe $u$ en los dos casos.
\end{ideaclave}

\subsection*{Lo que se usa de los capítulos anteriores}

Cuatro herramientas aparecen en todo el capítulo. Las funciones de comparación y la estabilidad uniforme se estudian en el capítulo~9; aquí se recuerda lo indispensable.

\begin{definicion}[label=def:av-comparacion]{Funciones de comparación}
$\alpha:[0,r)\to[0,\infty)$ es de \clave{clase $\K$} si es continua, estrictamente creciente y $\alpha(0)=0$; es de \clave{clase $\Kinf$} si además $r=\infty$ y $\alpha(s)\to\infty$ cuando $s\to\infty$. $\beta:[0,r)\times[0,\infty)\to[0,\infty)$ es de \clave{clase $\KL$} si $\beta(\cdot,s)\in\K$ para cada $s$ y $\beta(r,\cdot)$ es decreciente con $\beta(r,s)\to0$ cuando $s\to\infty$.
\end{definicion}

\begin{enumerate}
  \item[H1.] \textbf{Sándwich.} Si $V(t,x)$ es continua y $W_1(x)\le V(t,x)\le W_2(x)$ con $W_1,W_2$ definidas positivas en $\Bola r$, hay $\alpha_1,\alpha_2\in\K$ con $\alpha_1(\norm x)\le V(t,x)\le\alpha_2(\norm x)$; si $V$ es radialmente no acotada y el dominio es $\R^n$, se pueden tomar en $\Kinf$. Para $V=x\T Px$, $\lmin(P)\norm x^2\le V\le\lmax(P)\norm x^2$. Como $\alpha_1\le\alpha_2$,
  \begin{equation}
    \alpha_1^{-1}\bigl(\alpha_2(\mu)\bigr)\ge\mu\qquad\text{para todo }\mu>0 .
    \label{eq:av-a1a2}
  \end{equation}
  \item[H2.] \textbf{Conjuntos de nivel entre dos bolas fijas.} Con $\Nivel{t,c}=\{x:V(t,x)\le c\}$,
  \begin{equation}
    \Bola{\alpha_2^{-1}(c)}\subset\Nivel{t,c}\subset\Bola{\alpha_1^{-1}(c)}\qquad\text{para todo }t .
    \label{eq:av-inclusiones}
  \end{equation}
  El conjunto de nivel puede cambiar con $t$, pero queda encerrado entre dos bolas que no cambian: es la razón técnica de la uniformidad respecto a $t_0$.
  \item[H3.] \textbf{Comparación} (\cref{lem:pm-comparacion}). Si $\dot v\le-\alpha(v)$ con $\alpha\in\K$ localmente Lipschitz, entonces $v(t)\le\sigma(v(t_0),t-t_0)$ con $\sigma\in\KL$, la solución de $\dot y=-\alpha(y)$.
  \item[H4.] \textbf{Young.} Para $p,q>1$ con $1/p+1/q=1$ y $\lambda>0$,
  \begin{equation}
    \abs{yz}\le\frac{\lambda^p}{p}\abs y^p+\frac1{q\lambda^q}\abs z^q ;
    \label{eq:av-young}
  \end{equation}
  con $p=q=2$ y $\lambda=1$, $\abs{yz}\le\tfrac12y^2+\tfrac12z^2$.
\end{enumerate}

\begin{historia}[De la estabilidad al diseño robusto]
La idea de que la estabilidad implica la existencia de una función de Lyapunov es de Massera (1949 y 1956)~\citep{Mas56} y de Kurzweil; las versiones que se usan hoy están en \citet[\S4.7]{Kha02} y en \citet[cap.~2]{KeBr23}. Artstein mostró en 1983 que la existencia de una realimentación continua equivale a la de una función de control de Lyapunov~\citep{Ars83}, y Sontag dio en 1989 la fórmula explícita~\citep{Son89}. El mismo año, Sontag propuso la estabilidad entrada--estado~\citep{Son89a}, que con la caracterización de Sontag y Wang~\citep{SoWa95} y el teorema de pequeña ganancia de Jiang, Teel y Praly~\citep{JTP94} cambió la manera de diseñar controladores no lineales. Barbălat publicó su lema en 1959~\citep{Bar59}, en un artículo sobre sistemas de ecuaciones diferenciales no lineales.
\end{historia}

\begin{pausa}
Antes de seguir, prediga: (1) si $\Vdot<0$ en todo punto fuera de una bola $\Bola\mu$, ¿las trayectorias terminan dentro de $\Bola\mu$? (2) Si el origen de $\dot x=f(x)$ es estable, ¿existe siempre una $V(x)$, sin $t$, con $\Vdot\le0$? (3) Si $u=k(x)$ estabiliza globalmente, ¿una perturbación pequeña y acotada en $u$ deja el estado acotado? (4) Si $\dot x_2=f_2(x_2)$ y $\dot x_1=f_1(x_1,0)$ tienen el origen GAS, ¿la cascada $\dot x_1=f_1(x_1,x_2)$, $\dot x_2=f_2(x_2)$ es GAS?
\respuesta{Las cuatro respuestas son \emph{no}, por razones distintas. (1) Terminan en el menor conjunto de nivel de $V$ que contiene a $\Bola\mu$, que puede sobresalir de la bola (\cref{tr:av-cota}). (2) Hay sistemas invariantes en el tiempo estables sin ninguna $V(x)$ continua (\cref{ej:av-sinV}). (3) Una ley que cancela exactamente puede dejar escapar el estado en tiempo finito con un error de actuador arbitrariamente pequeño (\cref{ej:av-rediseno}). (4) $\dot x_1=-x_1+x_1^2x_2$, $\dot x_2=-x_2$ escapa en tiempo finito si $x_1(0)x_2(0)>2$ (\cref{ej:av-cascada}).}
\end{pausa}

%% ---------------------------------------------------------------------
\section{Acotamiento y acotamiento último uniforme}
\label{sec:av-uub}
%% ---------------------------------------------------------------------

Toda la teoría de estabilidad vista hasta ahora presupone un \emph{equilibrio}: $f(t,0)=0$ para todo $t$. En la práctica hay al menos tres situaciones en que ese punto no existe:
\begin{enumerate}
  \item \textbf{Perturbación persistente.} En $\dot x=f(x)+w(t)$ con $w\not\equiv0$, el origen deja de ser equilibrio aunque $f(0)=0$.
  \item \textbf{Error de modelado.} Lo que un controlador no cancela aparece como un término $\Delta(t,e)$ en la dinámica del error, que se sabe acotado pero no que se anule en $e=0$.
  \item \textbf{Control práctico.} Con ruido, cuantización o fricción seca, pedir $e(t)\to0$ es irreal; lo que se especifica es $\norm{e(t)}\le b$ después de un transitorio.
\end{enumerate}
La pregunta ya no es si $x(t)\to0$, sino dos cosas más modestas: ¿la solución queda \emph{acotada}, con una cota que no dependa de $t_0$? ¿Termina dentro de una \emph{bola fija}, después de un tiempo que tampoco depende de $t_0$?

\begin{ejemplo}[label=ej:av-escalar]{Un ejemplo que se puede resolver}
Estudie $\dot x=-x+\delta\sin t$ con $x(t_0)=a$ y $a>\delta>0$.
\solucion No hay equilibrio: $f(t,0)=\delta\sin t\neq0$. Pero la ecuación es lineal y, por variación de constantes,
\[
  x(t)=e^{-(t-t_0)}a+\delta\int_{t_0}^te^{-(t-\tau)}\sin\tau\,\dif\tau .
\]
\emph{Una cota uniforme.} Con $\abs{\sin\tau}\le1$ dentro de la integral,
\[
  \abs{x(t)}\le e^{-(t-t_0)}a+\delta\bigl(1-e^{-(t-t_0)}\bigr),
\]
una combinación convexa de $a$ y $\delta$, luego $\abs{x(t)}\le a$ para todo $t\ge t_0$, con una cota que no depende de $t_0$.
\emph{Una cota última.} Para cualquier $b$ con $\delta<b<a$, la combinación convexa queda por debajo de $b$ desde
\[
  t-t_0\ \ge\ T(a,b)=\ln\frac{a-\delta}{b-\delta} .
\]
$T$ depende del tamaño inicial y de la cota pedida, nunca de $t_0$. Cuando $b\downarrow\delta$, $T\to\infty$: se puede pedir una cota tan cercana a $\delta$ como se quiera, pero se paga con espera. Con $a=14$, $b=1.05$ y $\delta=1$, $T=\ln260=5.56$.
\emph{La cota real.} La solución particular es $\frac\delta2(\sin t-\cos t)=\frac{\delta}{\sqrt2}\sin(t-\frac\pi4)$, de amplitud $\delta/\sqrt2<\delta$. La precisión se perdió al reemplazar $\sin\tau$ por $\abs{\sin\tau}\le1$: la cota $\delta$ es la de la peor perturbación de amplitud $\delta$ (una constante $\pm\delta$), no la de una senoidal (\cref{fig:av-escalar}).
\end{ejemplo}

\begin{figure}[ht]
\centering
\includegraphics[width=\linewidth]{cap10/f_uub_escalar}
\caption{$\dot x=-x+\sin t$ desde $x(0)\in\{4,2,-4\}$. La envolvente (coral) es la cota $e^{-t}a+\delta(1-e^{-t})$ con $a=4$; desde $T(4,1.5)=\ln6=1.79$ las soluciones quedan en la banda $\abs x\le1.5$ (verde). Las punteadas marcan la amplitud real $\delta/\sqrt2$. A la derecha, el precio de afinar la cota: $T(a,b)$ crece sin cota cuando $b\to\delta$.}
\label{fig:av-escalar}
\end{figure}

\codigorepo{cap10_avanzada/C10_Ej1_UUBescalar.m}{escalar}{La solución cerrada, la envolvente y el tiempo de espera $T(a,b)$}

\video{V10_1_banda_uub}{Afinar la cota última cuesta tiempo}

Se considera $\dot x=f(t,x)$ con $f$ continua a trozos en $t$ y localmente Lipschitz en $x$, \emph{sin} suponer $f(t,0)=0$.

\begin{definicion}[label=def:av-uub]{UB y UUB \citep[Def.~4.6]{Kha02}}
Las soluciones son \clave{uniformemente acotadas (UB)} si existe $c>0$, independiente de $t_0\ge0$, y para cada $a\in(0,c)$ existe $\beta(a)>0$, independiente de $t_0$, tal que
\[
  \norm{x(t_0)}\le a\implica\norm{x(t)}\le\beta\qquad\forall\,t\ge t_0 .
\]
Son \clave{uniformemente últimamente acotadas (UUB) con cota última $b$} si existen $b,c>0$, independientes de $t_0$, y para cada $a\in(0,c)$ existe $T(a,b)\ge0$, independiente de $t_0$, tal que
\[
  \norm{x(t_0)}\le a\implica\norm{x(t)}\le b\qquad\forall\,t\ge t_0+T .
\]
Son \emph{globalmente} UB o UUB si las propiedades valen para $a$ arbitrariamente grande. Si el sistema es invariante en el tiempo se omite ``uniformemente''.
\end{definicion}

Hay que leer el papel de cada constante: $c$ delimita las condiciones iniciales permitidas; en UB la cota $\beta(a)$ crece con $a$ (quien empieza lejos puede pasearse lejos); en UUB la cota $b$ es fija y lo único que crece con $a$ es la espera $T(a,b)$. ``Uniforme'' significa lo mismo que en el capítulo~9: las constantes no dependen de $t_0$.

\begin{trampa}[label=tr:av-noestable]{``UUB es estabilidad con otra cota''}
\plausible{En la estabilidad hay una cota $\varepsilon$ y aquí una cota $b$: es la misma idea.}
En la estabilidad, para \emph{todo} $\varepsilon>0$ hay un $\delta(\varepsilon)$. En UUB la cota no se puede pedir tan pequeña como se quiera: hay un valor por debajo del cual ninguna $b$ funciona (en el \cref{ej:av-escalar}, ninguna $b<\delta/\sqrt2$ sirve para $w=\delta\sin t$, y ninguna $b\le\delta$ se puede garantizar solo con $\abs w\le\delta$).
\correcto{La cota última es un dato del problema, no un parámetro libre. Si las hipótesis del teorema de UUB valen para todo $\mu>0$, la cota tiende a cero con $\mu$ y se recupera la estabilidad asintótica uniforme: UUB es ``UAS hasta una bola''.}
\end{trampa}

\subsection{El argumento de los dos conjuntos de nivel}

Repetimos el \cref{ej:av-escalar} con $V(x)=\tfrac12x^2$, \emph{sin usar la solución}. A lo largo de las trayectorias,
\[
  \Vdot=-x^2+\delta x\sin t\le-x^2+\delta\abs x=-\abs x\bigl(\abs x-\delta\bigr),
\]
que no es negativa cerca del origen, pero sí fuera de la bola $\abs x\le\delta$.

\emph{Paso 1, acotamiento.} Para $c>\delta^2/2$, en la frontera de $\Nivel c=\{V\le c\}$ se tiene $\abs x=\sqrt{2c}>\delta$ y $\Vdot<0$. Una trayectoria que quisiera salir de $\Nivel c$ tendría que cruzar la frontera con $\Vdot\ge0$: $\Nivel c$ es \clave{positivamente invariante}. Con $c=a^2/2$, $\abs{x(t)}\le a$ para todo $t\ge t_0$.

\emph{Paso 2, acotamiento último.} Para $\delta^2/2<\varepsilon<c$, en la \clave{corona} $\Lambda=\{\varepsilon\le V\le c\}$ se tiene $\abs x\ge s=\sqrt{2\varepsilon}>\delta$ y $\Vdot\le-(s^2-\delta s)=-k<0$. Mientras la trayectoria esté en $\Lambda$, $V$ baja por lo menos a ritmo $k$, así que sale de la corona antes de
\begin{equation}
  T\le\frac{c-\varepsilon}{k}=\frac{c-\varepsilon}{2\varepsilon-\delta\sqrt{2\varepsilon}},
  \label{eq:av-Tlyap}
\end{equation}
y entra a $\Nivel\varepsilon$, que también es invariante. La cota última $\sqrt{2\varepsilon}$ puede acercarse a $\delta$ tanto como se quiera, y otra vez $T\to\infty$. No se usó la solución: solo $V$ y la cota de $\Vdot$. Esto es lo que se generaliza.

\begin{trampa}[label=tr:av-T]{``El argumento de Lyapunov da el tiempo de espera''}
\plausible{\eqref{eq:av-Tlyap} da un $T$ explícito, que debe ser el tiempo que tarda la solución en llegar.}
Da una cota superior, y muy pesimista. Con $a=14$, $\delta=1$ y $b=1.05$: la solución exacta llega en $\ln260=5.56$; \eqref{eq:av-Tlyap} da $c=98$, $\varepsilon=0.551$, $k=b^2-\delta b=0.0525$ y $T\le1856$. El argumento supone que $V$ baja siempre al ritmo mínimo $k$, que es el del borde interior de la corona.
\correcto{La garantía de Lyapunov es cualitativamente correcta y cuantitativamente conservadora. Sirve para todas las perturbaciones con $\abs w\le\delta$, no para una en particular.}
\end{trampa}

\subsection{El teorema y su geometría}

En $\R^n$ y con $V(t,x)$, las hipótesis son $\alpha_1(\norm x)\le V\le\alpha_2(\norm x)$ y $\Vdot\le-W_3(x)$ para $\norm x\ge\mu$. La región donde no se sabe nada de $\Vdot$ es la bola $\Bola\mu$, pero las trayectorias no quedan atrapadas en bolas sino en conjuntos de nivel de $V$. El menor que contiene a $\Bola\mu$ es $\Nivel{t,\eta}$ con $\eta=\alpha_2(\mu)$, y por \eqref{eq:av-inclusiones}
\begin{equation}
  \Bola\mu\ \subset\ \Nivel{t,\eta}\ \subset\ \Bola{\alpha_1^{-1}(\eta)} .
  \label{eq:av-cadena}
\end{equation}

\begin{teorema}[label=thm:av-uub]{Acotamiento último \citep[Teorema~4.18]{Kha02}}
Sea $D\subset\R^n$ un dominio que contiene al origen y $V:[0,\infty)\times D\to\R$ de clase $\Cont1$ con
\[
  \alpha_1(\norm x)\le V(t,x)\le\alpha_2(\norm x),\qquad \jac Vt+\jac Vxf(t,x)\le-W_3(x)\quad\forall\,\norm x\ge\mu>0,
\]
$\alpha_1,\alpha_2\in\K$ y $W_3$ continua y definida positiva. Sea $r>0$ con $\Bola r\subset D$ y $\mu<\alpha_2^{-1}(\alpha_1(r))$. Entonces existe $\beta\in\KL$ y, para cada condición inicial con $\norm{x(t_0)}\le\alpha_2^{-1}(\alpha_1(r))$, existe $T\ge0$ (que depende de $x(t_0)$ y de $\mu$, no de $t_0$) tal que
\begin{align}
  \norm{x(t)}&\le\beta\bigl(\norm{x(t_0)},t-t_0\bigr), && t_0\le t\le t_0+T, \label{eq:av-transitorio}\\
  \norm{x(t)}&\le\alpha_1^{-1}\bigl(\alpha_2(\mu)\bigr), && t\ge t_0+T. \label{eq:av-ultima}
\end{align}
Si $D=\R^n$ y $\alpha_1\in\Kinf$, valen para toda condición inicial y todo $\mu$.
\end{teorema}

\begin{demostracion}
(a) Con $c=\alpha_1(r)$, $\Nivel{t,c}\subset\Bola r\subset D$, y la condición inicial está en $\Nivel{t_0,c}$. (b) Con $\eta=\alpha_2(\mu)<c$, en la corona $\{\eta\le V\le c\}$ se tiene $\norm x\ge\mu$, luego $\Vdot\le-W_3(x)$. (c) En las fronteras $V=c$ y $V=\eta$, $\Vdot<0$: los dos conjuntos son positivamente invariantes y la solución existe para todo $t$. (d) $W_3$ tiene un mínimo $k>0$ en el compacto $\{\mu\le\norm x\le r\}$, así que la trayectoria entra a $\Nivel{t,\eta}$ antes de $T=(V(t_0,x(t_0))-\eta)/k$ y no sale; por \eqref{eq:av-cadena} vale \eqref{eq:av-ultima}. (e) Con $W_3\ge\alpha_3(\norm x)$, en la corona $\Vdot\le-\alpha_3(\alpha_2^{-1}(V))$, y el lema de comparación da $V\le\sigma(V(t_0),t-t_0)$, de donde $\beta(s,\tau)=\alpha_1^{-1}(\sigma(\alpha_2(s),\tau))$. Como las $\alpha_i$ no dependen de $t$, nada depende de $t_0$. (f) Si $\alpha_1\in\Kinf$, $r$ se puede tomar tan grande como haga falta.
\end{demostracion}

Cada hipótesis tiene su razón: $\mu<\alpha_2^{-1}(\alpha_1(r))$ garantiza que el conjunto trampa quepa donde $V$ está definida; la restricción sobre $x(t_0)$ pone la condición inicial en un conjunto de nivel dentro de $\Bola r$; que las $\alpha_i$ no dependan de $t$ da la uniformidad.

\begin{trampa}[label=tr:av-cota]{``La cota última es $\mu$''}
\plausible{Fuera de $\Bola\mu$ la función $V$ decrece, así que las trayectorias terminan en $\Bola\mu$.}
Solo se sabe que $V$ decrece \emph{fuera} de $\Bola\mu$; lo que no se abandona es un \emph{conjunto de nivel} de $V$, y el menor que contiene a $\Bola\mu$ llega hasta $\norm x=\alpha_1^{-1}(\alpha_2(\mu))\ge\mu$, por \eqref{eq:av-a1a2}. $\Bola\mu$ no tiene por qué ser invariante: dentro de ella no hay información sobre $\Vdot$.
\correcto{La cota garantizada es $b=\alpha_1^{-1}(\alpha_2(\mu))$, igual a $\mu$ solo si $V$ es ``redonda'' a esa escala ($\alpha_1(\mu)=\alpha_2(\mu)$, como con $V=\tfrac12\norm x^2$). El comportamiento real puede quedar por encima de $\mu$ o muy por debajo.}
\end{trampa}

\begin{ejemplo}[label=ej:av-oscilador]{La bola no es invariante}
Para el oscilador poco amortiguado $\dot x=Ax+Bw$, $A=\bigl[\begin{smallmatrix}0&1\\-4&-0.2\end{smallmatrix}\bigr]$, $B=\bigl[\begin{smallmatrix}0\\1\end{smallmatrix}\bigr]$, $\abs w\le0.2$, encuentre $\mu$ y la cota última con $V=x\T Px$, $A\T P+PA=-I$.
\solucion La ecuación de Lyapunov da $P=\bigl[\begin{smallmatrix}12.525&0.125\\0.125&3.125\end{smallmatrix}\bigr]$, con $\lmin=3.123$, $\lmax=12.527$ y $\norm{PB}=3.127$. Entonces $\Vdot=-\norm x^2+2x\T PBw\le-\norm x^2+2\norm{PB}\,0.2\norm x$ y, separando una fracción $\theta\in(0,1)$,
\[
  \Vdot\le-(1-\theta)\norm x^2\qquad\forall\,\norm x\ge\mu=\frac{2\norm{PB}\cdot0.2}{\theta}.
\]
Con $\theta=\tfrac12$, $\mu=2.50$ y $b=\mu\sqrt{\lmax/\lmin}=5.01$ (con $\theta=0.9$, $\mu=1.39$). Una trayectoria que nace en $(2.45,0)$, dentro de $\Bola\mu$, con $w=0.2\sin1.3t$, sube hasta $\norm x=4.50$ y nunca sale de $\Nivel\eta$ (\cref{fig:av-conjuntos}). La excursión ocurre \emph{incluso con $w\equiv0$}: el oscilador sigue una elipse casi conservativa de semiejes $2.45$ y $\approx4.9$. El problema es la excentricidad de los conjuntos de nivel ($\lmax/\lmin\approx4$), no la perturbación.
\end{ejemplo}

\begin{figure}[ht]
\centering
\includegraphics[width=\linewidth]{cap10/f_conjuntos}
\caption{Izquierda: el \cref{ej:av-oscilador}. $\Bola\mu$ (coral), $\Nivel\eta$ (azul) y $\Bola b$ (verde); la trayectoria (negra) nace dentro de $\Bola\mu$, sale de ella y nunca abandona $\Nivel\eta$. Derecha: el resorte endurecido del \cref{ej:av-resorte}, $\norm{x(t)}$ con $w=0.3\cos t$ frente a $\mu=0.745$ y la cota garantizada $b=1.30$.}
\label{fig:av-conjuntos}
\end{figure}

\codigorepo{cap10_avanzada/C10_Ej2_Conjuntos.m}{conjuntos}{$P$, $\mu$ y $b$ del oscilador, y la trayectoria que sale de $\Bola\mu$ sin salir de $\Nivel\eta$}

\video{V10_2_bola_no_invariante}{La bola no es invariante; el conjunto de nivel sí}

\begin{ejemplo}[label=ej:av-resorte]{Resorte endurecido forzado \citep[Ejemplo~4.23]{Kha02}}
Para $\dot x_1=x_2$, $\dot x_2=-(1+x_1^2)x_1-x_2+w(t)$, $\abs{w}\le M$, encuentre una cota última.
\solucion Con $w=0$ el origen es GAS. Se prueba $V=x\T Px+\tfrac12x_1^4$ (el término cuártico es la energía del resorte endurecido). Al derivar aparecen términos cruzados $x_1x_2$ y $x_1^3x_2$ sin signo; se anulan con $p_{22}=1$ y $p_{11}=p_{12}+1$, y con $p_{12}=\tfrac12$ queda $P=\bigl[\begin{smallmatrix}3/2&1/2\\1/2&1\end{smallmatrix}\bigr]$ y
\[
  \Vdot=-x_1^2-x_2^2-x_1^4+(x_1+2x_2)\,w\le-\norm x^2+\sqrt5\,M\norm x\le-(1-\theta)\norm x^2\quad\text{si }\norm x\ge\mu=\frac{\sqrt5M}{\theta}.
\]
Con $\lambda(P)=(5\mp\sqrt5)/4$, $\alpha_1=0.691\norm x^2$ y $\alpha_2=1.809\norm x^2+\tfrac12\norm x^4$: el resultado es global, con
\[
  b=\alpha_1^{-1}\bigl(\alpha_2(\mu)\bigr)=\sqrt{\frac{\lmax\mu^2+\mu^4/2}{\lmin}} .
\]
Con $M=0.3$ y $\theta=0.9$: $\mu=0.745$ y $b=1.30$. La simulación con $w=0.3\cos t$ da $\limsup\norm{x(t)}=0.30$ (con $w=0.3\cos2t$, $0.17$; con $w\equiv0.3$, $0.28$). La cadena $0.30<\mu<b$ resume las dos trampas: $\mu$ no es lo que se garantiza, y lo que se garantiza es conservador.
\end{ejemplo}

\codigorepo{cap10_avanzada/C10_Ej2_Conjuntos.m}{resorte}{El resorte endurecido: $\mu$, $b$ y el $\limsup$ simulado para tres perturbaciones}

Dos corolarios se usan a menudo. Con las mismas cotas sobre $V$: si $\Vdot\le0$ para $\norm x\ge\mu$, las soluciones que nacen en $\Nivel{t_0,\eta}$, $\eta\ge\alpha_2(\mu)$, cumplen $\norm{x(t)}\le\alpha_1^{-1}(\eta)$ (UB); si $\Vdot\le-W(x)$ con $W>0$ para $\norm x>\mu$, son UUB con $\limsup\norm{x(t)}\le\alpha_1^{-1}(\alpha_2(\mu))$. En los dos hay que decir \emph{para qué condiciones iniciales} valen; en el caso local, sin esa restricción, son falsos.

\begin{pausa}
Para $\dot x=\delta\sin t$, ¿son las soluciones UB? ¿UUB? ¿Qué hipótesis del \cref{thm:av-uub} falla?
\respuesta{$x(t)=x(t_0)+\delta(\cos t_0-\cos t)$, así que $\abs{x(t)}\le\abs{x(t_0)}+2\delta$: son UB. No son UUB: una solución que empieza lejos se queda lejos. Con $V=\tfrac12x^2$, $\Vdot=\delta x\sin t$ no tiene signo en ninguna región, y no existe $W_3$ con $\Vdot\le-W_3$ fuera de una bola: el sistema no tiene disipación.}
\end{pausa}

%% ---------------------------------------------------------------------
\section{Teoremas conversos de Lyapunov}
\label{sec:av-conversos}
%% ---------------------------------------------------------------------

El teorema directo de Lyapunov dice: \emph{si existe $V$, hay estabilidad}. Un teorema converso dice lo contrario: \emph{si hay estabilidad de cierto tipo, existe una $V$ con ciertas propiedades}. Tiene tres usos, en este orden de importancia práctica:
\begin{enumerate}
  \item \textbf{No buscar en vano.} Si el origen es UAS, alguna $V$ existe: si no la encuentro, el problema es mi candidata, no el sistema.
  \item \textbf{Robustez.} La $V$ que garantiza el converso tiene propiedades \emph{cuantitativas}, como $\norm{\partial V/\partial x}\le c_4\norm x$, que permiten probar que la estabilidad sobrevive a perturbaciones (\cref{lem:av-robustez}).
  \item \textbf{Diseño modular.} En CLF, \emph{backstepping} o ISS se supone ``un subsistema ya estabilizado, con su función de Lyapunov''. El converso garantiza que esa función existe; el rediseño ISS de la \S\ref{sec:av-iss} lo invoca explícitamente.
\end{enumerate}

\begin{trampa}[label=tr:av-converso]{``El converso me da la $V$''}
\plausible{Si el teorema garantiza que $V$ existe, su demostración me dice cómo calcularla.}
Las demostraciones construyen $V$ \emph{integrando las soluciones del sistema}, que es justamente lo que se quería evitar al usar Lyapunov. Salvo en el caso lineal, la $V$ existe pero no se puede escribir. \citet[\S2.5]{KeBr23} lo dicen sin rodeos: la prueba no ofrece un punto de partida para construir una función de Lyapunov.
\correcto{El converso sirve para razonar (qué propiedades tiene $V$, qué robustez implican), no para calcular. Para calcular siguen haciendo falta ingenio, energía física, el gradiente variable o la ecuación de Lyapunov de la linealización.}
\end{trampa}

\subsection{La energía futura}

El caso lineal muestra la idea de todas las pruebas: \clave{$V(x)$ es la energía futura total de la solución que parte de $x$}.

\begin{ejemplo}[label=ej:av-energia]{La ecuación de Lyapunov aparece sola}
Para $\dot x=Ax$ con $A$ Hurwitz, sea $\sol(\tau;x)=e^{A\tau}x$ la solución que parte de $x$ y
\[
  V(x)=\int_0^\infty\norm{\sol(\tau;x)}^2\dif\tau=x\T\Bigl(\int_0^\infty e^{A\T\tau}e^{A\tau}\dif\tau\Bigr)x=x\T Px .
\]
Muestre que $\Vdot=-\norm x^2$ y calcule $P$ para $A=\bigl[\begin{smallmatrix}0&1\\-1&-1\end{smallmatrix}\bigr]$.
\solucion La integral converge porque $\norm{e^{A\tau}}\le ke^{-\lambda\tau}$. Por la propiedad de semigrupo, $\sol(\tau;x(t))=x(t+\tau)$, y a lo largo de una trayectoria
\[
  V\bigl(x(t)\bigr)=\int_t^\infty\norm{x(s)}^2\dif s\quad\implica\quad\frac{\dif}{\dif t}V\bigl(x(t)\bigr)=-\norm{x(t)}^2 :
\]
al avanzar en el tiempo, la energía futura solo pierde el presente. Como también $\frac{\dif}{\dif t}x\T Px=x\T(A\T P+PA)x$, igualando para todo $x$ resulta $A\T P+PA=-I$: la ecuación de Lyapunov del \cref{thm:ly-ecuacion} es la condición de que la energía futura se pierda al ritmo $\norm x^2$. Con $P=\bigl[\begin{smallmatrix}p&q\\q&r\end{smallmatrix}\bigr]$ la ecuación da $-2q=-1$, $p-q-r=0$, $2(q-r)=-1$, de donde
\[
  P=\begin{bmatrix}1.5&0.5\\0.5&1\end{bmatrix},\qquad V(1,0.5)=2.25 ,
\]
que es exactamente el área bajo $\norm{\sol(\tau;x)}^2$ desde $x=(1,0.5)$ (\cref{fig:av-energia}).
\end{ejemplo}

\begin{figure}[ht]
\centering
\includegraphics[width=\linewidth]{cap10/f_energia_futura}
\caption{La función de Lyapunov como energía futura. Izquierda: $\norm{\sol(\tau;x_0)}^2$ desde $x_0=(1,0.5)$; el área sombreada vale $V(x_0)=x_0\T Px_0=2.25$. Derecha: $V(x(t))=\int_t^\infty\norm{x}^2$ sobre la trayectoria; su pendiente es $-\norm{x(t)}^2$.}
\label{fig:av-energia}
\end{figure}

\codigorepo{cap10_avanzada/C10_Ej3_EnergiaFutura.m}{energia}{La integral de la energía futura contra $x\T Px$ con \cmd{lyap}}

\video{V10_3_energia_futura}{La energía futura se gasta en el presente}

\subsection{El converso exponencial}

\begin{teorema}[label=thm:av-converso]{Converso exponencial \citep[Teorema~4.14]{Kha02}}
Sea $f:[0,\infty)\times D\to\R^n$ de clase $\Cont1$, $D=\Bola r$, con $\norm{\partial f/\partial x}\le L$ en $D$ uniformemente en $t$ y $f(t,0)=0$. Sean $k,\lambda,r_0>0$ con $r_0<r/k$. Si
\[
  \norm{x(t)}\le k\norm{x(t_0)}e^{-\lambda(t-t_0)}\qquad\forall\,x(t_0)\in\Bola{r_0},\ t\ge t_0\ge0,
\]
existen $V:[0,\infty)\times\Bola{r_0}\to\R$ y constantes $c_1,\dots,c_4>0$ tales que
\begin{align}
  &c_1\norm x^2\le V(t,x)\le c_2\norm x^2, \label{eq:av-conv1}\\
  &\jac Vt+\jac Vxf(t,x)\le-c_3\norm x^2, \label{eq:av-conv2}\\
  &\Bigl\|\jac Vx\Bigr\|\le c_4\norm x . \label{eq:av-conv3}
\end{align}
Si $r=\infty$ y el origen es GES, todo vale en $\R^n$. Si $f$ no depende de $t$, $V$ tampoco.
\end{teorema}

La prueba toma la energía futura en una \emph{ventana finita} de longitud $h>0$,
\begin{equation}
  V(t,x)=\int_t^{t+h}\sol\T(\tau;t,x)\,\sol(\tau;t,x)\,\dif\tau ,
  \label{eq:av-Vconverso}
\end{equation}
y cada constante sale de una hipótesis distinta:
\begin{itemize}
  \item $c_2$, de la cota exponencial: $V\le k^2\norm x^2\int_0^he^{-2\lambda s}\dif s$, así que $c_2=k^2(1-e^{-2\lambda h})/(2\lambda)$.
  \item $c_1$, de la cota del jacobiano: con $\norm{f}\le L\norm x$, la solución \emph{no decae demasiado rápido}, $\norm{\sol(\tau)}\ge\norm xe^{-L(\tau-t)}$ (Gronwall, \cref{lem:pm-gronwall}), y $c_1=(1-e^{-2Lh})/(2L)$.
  \item $c_3$, de que la ventana es finita: $\frac{\dif}{\dif t}V=\norm{x(t+h)}^2-\norm{x(t)}^2\le-(1-k^2e^{-2\lambda h})\norm{x(t)}^2$; con $h=\ln(2k^2)/(2\lambda)$, $c_3=\tfrac12$.
  \item $c_4$, la novedad: la sensibilidad $\partial\sol/\partial x$ cumple la ecuación variacional y $\norm{\partial\sol/\partial x}\le e^{L(\tau-t)}$, así que $\norm{\partial V/\partial x}\le2k\norm x\int_0^he^{(L-\lambda)s}\dif s=c_4\norm x$. Con ventana infinita esta integral divergería si $L\ge\lambda$: por eso la ventana es finita.
\end{itemize}

\begin{table}[ht]
\centering\small
\begin{tabular}{@{}p{3.4cm}p{6.6cm}p{2.6cm}@{}}
\toprule
Hipótesis sobre el origen & Se garantiza una $V$ con\ldots & Referencia\\
\midrule
ES (local), $f\in\Cont1$, jacobiano acotado & $c_1\norm x^2\le V\le c_2\norm x^2$, $\Vdot\le-c_3\norm x^2$, $\norm{\partial V/\partial x}\le c_4\norm x$ en $\Bola{r_0}$ & \cite[Teo.~4.14]{Kha02}\\
GES, invariante en el tiempo & lo mismo, en $\R^n$, con $V(x)$ & \cite[Teo.~2.32]{KeBr23}\\
UAS (local), $f\in\Cont1$, jacobiano acotado & $\alpha_1\le V\le\alpha_2$, $\Vdot\le-\alpha_3(\norm x)$, $\norm{\partial V/\partial x}\le\alpha_4(\norm x)$, $\alpha_i\in\K$ & \cite{Mas56}; \cite[Teo.~4.16]{Kha02}\\
UGAS & $V$ suave, $\alpha_1,\alpha_2\in\Kinf$, $\Vdot\le-\rho(\norm x)$; $V(x)$ si $f$ es invariante en el tiempo & \cite[Teo.~2.31]{KeBr23}\\
Solo estable & $V(t,x)$ definida positiva con $\Vdot\le0$; en general \textbf{no} se puede quitar la $t$ & \cite[\S2.5.1]{KeBr23}\\
\bottomrule
\end{tabular}
\caption{Qué entrega cada teorema converso.}
\label{tab:av-conversos}
\end{table}

La fila UAS dice más que ``$\Vdot<0$'': da una cota \emph{uniforme} $\Vdot\le-\alpha_3(\norm x)$ del ritmo de decrecimiento. Esa uniformidad es la que da robustez frente a perturbaciones persistentes: sin ella no se podría aplicar el \cref{thm:av-uub} al sistema perturbado.

\begin{ejemplo}[label=ej:av-sinV]{Estable, invariante en el tiempo y sin $V(x)$ \citep[ec.~(2.51)]{KeBr23}}
Para
\[
  \dot x_1=x_2,\qquad\dot x_2=\sin^2\!\Bigl(\frac{\pi}{x_1^2+x_2^2}\Bigr)x_2-x_1 ,
\]
con $f(0)=0$, muestre que el origen es estable y que no existe una $V(x)$ continua, definida positiva, con $\Vdot\le0$.
\solucion \emph{Órbitas periódicas.} Sobre los círculos $\Gamma_n=\{\norm x^2=1/n\}$ el seno se anula y el campo es $(x_2,-x_1)$, tangente al círculo: cada $\Gamma_n$ es una órbita periódica (radios $1$, $1/\sqrt2=0.707$, $1/\sqrt3=0.577$,\ldots).
\emph{Entre órbitas la norma crece:} $\frac{\dif}{\dif t}\tfrac12\norm x^2=\sin^2(\pi/\norm x^2)\,x_2^2\ge0$, y las trayectorias espiralan hacia la órbita exterior.
\emph{Estable, no AS.} Las $\Gamma_n$ son barreras: dado $\varepsilon$, se toma $1/\sqrt n<\varepsilon$ y $\delta=1/\sqrt{n+1}$. Hay órbitas periódicas arbitrariamente cerca del origen, así que no es AS.
\emph{No hay $V(x)$.} Si existiera: (i) $V$ es constante sobre cada $\Gamma_n$ (si bajara en un tramo, al dar la vuelta $V(x(T))<V(x(0))=V(x(T))$); (ii) una espiral va de $\Gamma_{n+1}$ (en el pasado) a $\Gamma_n$ (en el futuro) sin que $V$ crezca, luego $V|_{\Gamma_n}\le V|_{\Gamma_{n+1}}$; (iii) encadenando y por continuidad, $V|_{\Gamma_1}\le V|_{\Gamma_2}\le\cdots\to V(0)=0$, que contradice $V>0$.
\end{ejemplo}

\begin{figure}[ht]
\centering
\includegraphics[width=0.75\linewidth]{cap10/f_sinV}
\caption{El \cref{ej:av-sinV}: órbitas periódicas $\Gamma_1,\Gamma_2,\Gamma_3$ (azul) y una trayectoria que espirala de $\Gamma_3$ hacia $\Gamma_2$. El origen es estable, pero ninguna $V(x)$ continua decrece a lo largo de todas las trayectorias.}
\label{fig:av-sinV}
\end{figure}

\codigorepo{cap10_avanzada/C10_Ej4_EstableSinV.m}{sinV}{Las órbitas $\Gamma_n$ y las espirales que las conectan}

Para sistemas invariantes en el tiempo, ``AS implica que existe $V(x)$'' es el teorema de Massera, pero ``estable implica que existe $V(x)$'' es falso: el converso para estabilidad simple necesita, en general, una $V(t,x)$ que dependa del tiempo \emph{aunque el sistema no dependa de él}. (El campo no es Lipschitz en el origen, pero las soluciones siguen siendo únicas: $\frac{\dif}{\dif t}\norm x^2\le2\norm x^2$ y la solución que nace en $0$ se queda en $0$.)

\subsection{Robustez frente a perturbaciones que se anulan}

\begin{lema}[label=lem:av-robustez]{Perturbaciones que se anulan \citep[Lema~9.1]{Kha02}}
Si el origen de $\dot x=f(t,x)$ es exponencialmente estable, $V$ es la función del \cref{thm:av-converso} y la perturbación cumple
\[
  \norm{g_p(t,x)}\le\nu\norm x\qquad\text{con}\qquad\nu<\frac{c_3}{c_4},
\]
el origen de $\dot x=f(t,x)+g_p(t,x)$ es exponencialmente estable (GES si todo vale globalmente).
\end{lema}

\begin{demostracion}
A lo largo del sistema perturbado,
$\Vdot\le-c_3\norm x^2+\bigl\|\jac Vx\bigr\|\norm{g_p}\le-(c_3-c_4\nu)\norm x^2\le-\frac{c_3-c_4\nu}{c_2}V$. Sin \eqref{eq:av-conv3} no habría forma de acotar el término cruzado $\jac Vxg_p$, porque de $g_p$ solo se conoce el tamaño.
\end{demostracion}

\begin{ejemplo}[label=ej:av-robustez]{El margen garantizado no es el margen real}
Para $A=\bigl[\begin{smallmatrix}0&1\\-1&-1\end{smallmatrix}\bigr]$ y la $P$ del \cref{ej:av-energia}, ¿qué perturbaciones tolera el sistema?
\solucion Con $V=x\T Px$: $c_3=1$ y $\partial V/\partial x=2x\T P$, así que $c_4=2\lmax(P)=2\cdot1.809=3.62$. El \cref{lem:av-robustez} garantiza estabilidad para $\nu<1/3.62=0.276$. Para la perturbación particular $g_p=\nu\,[0;\ x_1]$ el sistema es lineal con matriz $\bigl[\begin{smallmatrix}0&1\\-1+\nu&-1\end{smallmatrix}\bigr]$, Hurwitz si y solo si $\nu<1$. La garantía es menor porque cubre \emph{toda} $g_p$ con $\norm{g_p}\le\nu\norm x$, incluidas perturbaciones no lineales y variantes en el tiempo, mucho más dañinas que esta. La simulación con $\nu=0.27$, $0.9$ y $1.05$ da convergencia, convergencia lenta y divergencia.
\end{ejemplo}

\codigorepo{cap10_avanzada/C10_Ej5_Robustez.m}{robustez}{El margen $\nu^*=c_3/c_4$ contra el margen real de una perturbación particular}

\begin{pausa}
Para $\dot x=-x^3$, ¿existe $V$ con $c_1x^2\le V\le c_2x^2$ y $\Vdot\le-c_3x^2$? ¿Qué falla si se intenta la energía futura $\int_0^\infty\sol^2\dif\tau$?
\respuesta{No existe: si existiera, $\Vdot\le-(c_3/c_2)V$ y el origen sería exponencialmente estable, pero $\abs{x(t)}=\abs{x_0}/\sqrt{1+2x_0^2t}$ decae como $1/\sqrt t$. La energía futura diverge: $\int_0^\infty x^2/(1+2x^2\tau)\dif\tau=\tfrac12\ln(1+2x^2\tau)\big|_0^\infty=\infty$. No contradice a Massera: su construcción pondera la integral con una función $G(\norm\sol)$ que se aplana cerca de cero lo suficiente para que converja.}
\end{pausa}

%% ---------------------------------------------------------------------
\section{Funciones de control de Lyapunov y fórmula de Sontag}
\label{sec:av-clf}
%% ---------------------------------------------------------------------

En el análisis, $f$ está dada y se busca una $V$ con $\Vdot<0$. En el diseño se invierte el orden: \alerta{primero se escoge la ``energía'' $V$ y después la entrada que la hace bajar en cada punto}. Se considera un sistema invariante en el tiempo con entrada,
\[
  \dot x=f(x,u),\qquad x\in\R^n,\ u\in\R^m,\qquad f(0,0)=0,
\]
y se busca $u=k(x)$ que haga al origen de $\dot x=f(x,k(x))$ GAS imponiendo $\jac Vx f(x,k(x))\le-W(x)$ con $W$ definida positiva. Si en algún $x\neq0$ \emph{ninguna} $u$ logra $\Vdot<0$, esa $V$ no sirve para diseñar.

\begin{definicion}[label=def:av-clf]{Función de control de Lyapunov \citep{Ars83,Son83}}
$V:\R^n\to[0,\infty)$ de clase $\Cont1$, definida positiva y radialmente no acotada, es una \clave{función de control de Lyapunov (CLF)} para $\dot x=f(x,u)$ si
\[
  \inf_u\Bigl\{\jac Vx(x)\,f(x,u)\Bigr\}<0\qquad\forall\,x\neq0 .
\]
\end{definicion}

\begin{trampa}[label=tr:av-noclf]{``Si el sistema es estabilizable, cualquier $V$ cuadrática es CLF''}
\plausible{El doble integrador se estabiliza con $u=-x_1-x_2$; la energía $V=\tfrac12(x_1^2+x_2^2)$ debería servir para diseñar.}
Con $\dot x_1=x_2$, $\dot x_2=u$, $\Vdot=x_1x_2+x_2u=x_2(x_1+u)$. En los puntos $x_2=0$, $x_1\neq0$, $\Vdot=0$ \emph{para toda} $u$: $V$ no es CLF.
\correcto{La culpa es de $V$, no del sistema. Con $V=x\T Px$, $P=\bigl[\begin{smallmatrix}p&q\\q&r\end{smallmatrix}\bigr]>0$, en el conjunto $\Lie gV=0$ (es decir, $x_2=-qx_1/r$) se tiene $\Lie fV=-2q(pr-q^2)x_1^2/r^2$, negativa para $x_1\neq0$ si y solo si $q>0$: cualquier $P>0$ con término cruzado positivo es CLF.}
\end{trampa}

\subsection{El caso afín en el control}

Para $\dot x=f(x)+g(x)u$ con $f(0)=0$, la derivada de $V$ es afín en $u$:
\[
  \Vdot=\underbrace{\jac Vxf(x)}_{\ell_f(x)=\Lie fV}+\underbrace{\jac Vxg(x)}_{\ell_g(x)=\Lie gV}u .
\]
Si $\ell_g(x)\neq0$, el ínfimo sobre $u$ es $-\infty$. El único riesgo está donde $\ell_g(x)=0$, y allí $\Vdot=\ell_f(x)$ no depende de $u$. Por lo tanto, $V$ es CLF si y solo si
\begin{equation}
  \ell_g(x)=0\ \implica\ \ell_f(x)<0\qquad\forall\,x\neq0 .
  \label{eq:av-clfafin}
\end{equation}
Geométricamente, $\ell_g=0$ significa $\nabla V\perp g$: el control empuja \emph{tangente} al conjunto de nivel y no puede bajar $V$; ahí la deriva tiene que hacer el trabajo sola.

\begin{teorema}[label=thm:av-artstein]{Artstein \citep{Ars83}; versión global y fórmula: \citet{Son89}}
Sea $\dot x=f(x)+g(x)u$ con $f,g$ suaves y $f(0)=0$. Existe una realimentación $k$ \emph{continua en $\R^n$ y suave en $\R^n\setminus\{0\}$} que hace al origen GAS si y solo si existe $V$ suave, definida positiva y radialmente no acotada con
\begin{enumerate}
  \item (CLF) la condición \eqref{eq:av-clfafin};
  \item (\clave{propiedad de control pequeño}) para todo $\varepsilon>0$ existe $\delta>0$ tal que, si $0<\norm x<\delta$, hay un $u$ con $\norm u<\varepsilon$ y $\ell_f(x)+\ell_g(x)u<0$.
\end{enumerate}
\end{teorema}

La dirección ``$\Leftarrow$'' es constructiva: la da la fórmula de Sontag. La dirección ``$\Rightarrow$'' aplica un teorema converso al lazo cerrado (uno de los usos modulares de la \S\ref{sec:av-conversos}). La propiedad de control pequeño es exactamente lo que da la continuidad en el origen.

\begin{proposicion}[label=prop:av-sontag]{Fórmula universal de Sontag \citep{Son89}}
Si $V$ es una CLF para $\dot x=f(x)+g(x)u$, $u\in\R$, la realimentación
\begin{equation}
  k(x)=\begin{cases}-\dfrac{\ell_f(x)+\sqrt{\ell_f(x)^2+\ell_g(x)^4}}{\ell_g(x)}, & \ell_g(x)\neq0,\\[8pt] 0, & \ell_g(x)=0,\end{cases}
  \label{eq:av-sontag}
\end{equation}
da $\Vdot=-W(x)$ con $W=\sqrt{\ell_f^2+\ell_g^4}>0$ para todo $x\neq0$, es suave en $\R^n\setminus\{0\}$ y es continua en el origen si y solo si $V$ tiene la propiedad de control pequeño.
\end{proposicion}

\begin{demostracion}
\emph{Decrecimiento.} Si $\ell_g\neq0$, $\Vdot=\ell_f-\bigl(\ell_f+\sqrt{\ell_f^2+\ell_g^4}\bigr)=-\sqrt{\ell_f^2+\ell_g^4}$. Si $\ell_g=0$, $k=0$ y $\Vdot=\ell_f=-\abs{\ell_f}$ por \eqref{eq:av-clfafin}.
\emph{Suavidad lejos del origen.} Se escribe $k=-\varphi(\ell_f,\ell_g^2)\,\ell_g$ con $\varphi(a,p)=(a+\sqrt{a^2+p^2})/p$ y $\varphi(a,0)=0$ si $a<0$. $\varphi$ es la raíz de $p\varphi^2-2a\varphi-p=0$, cuya derivada respecto a $\varphi$ vale $2\sqrt{a^2+p^2}>0$ en $S=\{p>0\text{ o }a<0\}$; por la función implícita, $\varphi$ es analítica en $S$, y la condición CLF dice que $(\ell_f,\ell_g^2)\in S$ para $x\neq0$.
\emph{Continuidad en el origen.} Si $k$ es continua, $u=k(x)$ da la propiedad de control pequeño. Recíprocamente, para $0<\norm x<\delta$ con $\ell_g\neq0$: si $\ell_f\ge0$ hay $\abs u<\varepsilon$ con $\ell_f<\abs{\ell_g}\varepsilon$, y $\abs k\le(2\ell_f+\ell_g^2)/\abs{\ell_g}<2\varepsilon+\abs{\ell_g}$; si $\ell_f<0$, $\abs k\le\ell_g^2/\abs{\ell_g}=\abs{\ell_g}$. Como $\nabla V(0)=0$, $\ell_g(x)\to0$ y $\limsup_{x\to0}\abs{k(x)}\le2\varepsilon$ para todo $\varepsilon$.
\end{demostracion}

\begin{trampa}[label=tr:av-cuarta]{``Basta con $\ell_g^2$ dentro de la raíz''}
\plausible{La fórmula funcionaría igual con $\sqrt{\ell_f^2+\ell_g^2}$: también da $\Vdot<0$.}
Da $\Vdot<0$, pero no una ley continua. En un punto con $\ell_f=0$ y $\ell_g\to0$, la variante da $k\approx-\abs{\ell_g}/\ell_g=-\sgn(\ell_g)$, que salta de $-1$ a $+1$; con $\ell_g^4$, $k\approx-\ell_g^2/\ell_g=-\ell_g\to0$.
\correcto{El exponente 4 hace que dentro de la raíz aparezca $p^2=(\ell_g^2)^2$, con $p=\ell_g^2$ suave. Con $\ell_g^2$ aparecería $p=\abs{\ell_g}$, que no es suave donde $\ell_g=0$.}
\end{trampa}

Con varias entradas, $\ell_g$ es un vector fila y la fórmula se escribe \citep[ec.~(9.10)]{KeBr23}
\begin{equation}
  k(x)=-\biggl(\kappa+\frac{\ell_f+\sqrt{\ell_f^2+\norm{\ell_g}^4}}{\norm{\ell_g}^2}\biggr)\ell_g\T,\qquad\kappa\ge0,
  \label{eq:av-sontagvec}
\end{equation}
con $k=0$ donde $\ell_g=0$. El término $-\kappa\ell_g\T$ agrega $-\kappa\norm{\ell_g}^2$ a $\Vdot$ y, como se verá en la \S\ref{sec:av-iss}, da robustez frente a errores aditivos en el actuador.

\subsection{Anatomía: ¿cuánto control gasta?}

La fórmula se lee en tres regímenes. Donde $\ell_g=0$, $k=0$: la deriva baja $V$ sola. Si $\ell_f<0$ (la deriva ayuda) y $\abs{\ell_g}$ es pequeño, $k\approx-\ell_g^3/(2\abs{\ell_f})$: el control casi no actúa. Si $\ell_f>0$ (la deriva estorba), $k\approx-2\ell_f/\ell_g$: cancela la deriva y la invierte. \alerta{Sontag no gasta control donde la planta ya hace el trabajo}; la linealización exacta del \cref{cap:linealizacion} cancela toda la deriva, útil o no.

\begin{ejemplo}[label=ej:av-sontaglineal]{Planta lineal escalar}
Aplique la fórmula a $\dot x=f_0x+g_0u$ con $V=\tfrac12x^2$.
\solucion $\ell_f=f_0x^2$, $\ell_g=g_0x$ y
\[
  k(x)=-\frac{f_0+\sqrt{f_0^2+g_0^4}}{g_0}\,x,\qquad\dot x=-\sqrt{f_0^2+g_0^4}\,x .
\]
Con $f_0=g_0=1$, $k=-(1+\sqrt2)x=-2.414x$ y, desde $x(0)=2$, $u(0)=-4.83$. (Sin control, $u\equiv1$ daría $x(10)=3e^{10}-1=6.6\times10^4$.) Si $g_0\to0$ con $f_0>0$, la ganancia $\approx-2f_0/g_0\to-\infty$: un actuador débil debe empujar fuerte. Si $f_0<0$, la ganancia $\approx-g_0^3/(2\abs{f_0})\to0$: el control no se molesta.
\end{ejemplo}

\begin{trampa}[label=tr:av-escala]{``Universal significa sin perillas''}
\plausible{La fórmula no tiene parámetros: dada la CLF, la ley está determinada.}
Con $\tilde V=cV$, $\tilde\ell_f=c\ell_f$, $\tilde\ell_g=c\ell_g$ y $\tilde k=-(\ell_f+\sqrt{\ell_f^2+c^2\ell_g^4})/\ell_g\neq k$: $\ell_f^2$ y $\ell_g^4$ no tienen las mismas unidades.
\correcto{La escala de $V$ (o de $u$) es un parámetro de sintonía. ``Universal'' significa que sirve para cualquier CLF, no que no tenga perillas.}
\end{trampa}

\begin{ejemplo}[label=ej:av-sontagcos]{$\dot x=\cos x-x^3+u$}
Diseñe con Sontag y compare con la linealización exacta desde $x(0)=2$.
\solucion \emph{Primero:} $f(0)=\cos0=1\neq0$. Con $u=0$ el origen no es equilibrio, y la fórmula presupone $f(0)=0$; aplicada a ciegas con $V=\tfrac12x^2$ da una ley discontinua. Hay que separar lo que no se anula. Con $u=-\cos x+u_s$ (se cancela $\cos x$), $\dot x=-x^3+u_s$, $\ell_f=-x^4$, $\ell_g=x$ y
\[
  u_s=x^3-x\sqrt{x^4+1},\qquad\dot x=-x\sqrt{x^4+1} .
\]
(Otra opción legítima: $u=-1+v$, que solo compensa el valor $\cos0$, y Sontag sobre $\dot x=\cos x-1-x^3+v$.) La linealización exacta $u_{FL}=-\cos x+x^3-x$ da $\dot x=-x$. Desde $x(0)=2$, $u_{FL}(0)=6.42$ y $u_S(0)=-\cos2+8-2\sqrt{17}=0.17$: $u_{FL}$ cancela el término $-x^3$, que ayudaba; Sontag casi lo conserva. Y $u(t)\to-\cos0=-1$, no a cero (\cref{fig:av-sontag}).
\end{ejemplo}

\begin{trampa}[label=tr:av-regimen]{``Si el estado llega a cero, el control también''}
\plausible{En el equilibrio no hay nada que corregir, así que $u\to0$.}
En el \cref{ej:av-sontagcos}, $u(t)\to-1$: mantener $x=0$ exige compensar permanentemente la ``fuerza'' $\cos0=1$.
\correcto{Si $f(0)\neq0$, el control en régimen es el $u^*$ que hace del origen un equilibrio, $f(0)+g(0)u^*=0$. (En la gráfica de las diapositivas anteriores $u(t)$ arrancaba en $\approx-20$ y tendía a $0$: el diagrama de \simulink{} no implementaba la fórmula.)}
\end{trampa}

\begin{figure}[ht]
\centering
\includegraphics[width=\linewidth]{cap10/f_sontag}
\caption{Izquierda y centro: el \cref{ej:av-sontagcos}, $x(t)$ y $u(t)$ con Sontag (azul) y con linealización exacta (coral) desde $x(0)=2$; Sontag arranca con $u=0.17$ y la linealización con $6.42$; ambas tienden a $u^*=-1$. Derecha: el \cref{ej:av-controlpequeno}, $\abs{u(t)}$ crece sin cota mientras $x(t)\to0$.}
\label{fig:av-sontag}
\end{figure}

\codigorepo{cap10_avanzada/C10_Ej6_SontagLineal.m}{lineal}{La fórmula de Sontag en la planta escalar}
\codigorepo{cap10_avanzada/C10_Ej7_SontagCos.m}{sontag}{Sontag contra linealización exacta en $\dot x=\cos x-x^3+u$}

\begin{modelo}[RedisenoISS.slx]
\cmd{crearModelo\_RedisenoISS.m} arma el lazo $\dot x=x+(x^2+1)(k(x)+w)$ con un selector entre tres leyes (cancelación exacta, cancelación con amortiguamiento $-\Lie gV$ y Sontag con $\kappa$) y un escalón de perturbación $w$ en el actuador. Se usa en este capítulo dos veces: aquí para comparar el esfuerzo de control y en la \S\ref{sec:av-iss} para ver cuál ley sobrevive a la perturbación.
\end{modelo}

\begin{ejemplo}[label=ej:av-controlpequeno]{Cuando falla la propiedad de control pequeño}
Estudie $\dot x=x+x^2u$ con $V=\tfrac12x^2$.
\solucion $\ell_f=x^2$, $\ell_g=x^3$ se anula solo en $0$: $V$ es CLF, y
\[
  k(x)=-\frac{1+\sqrt{1+x^8}}{x},\qquad\dot x=-x\sqrt{1+x^8} .
\]
El origen del lazo cerrado es GAS, pero $\abs{u(t)}\to\infty$ mientras $x(t)\to0$: desde $x(0)=1$, $x(8)\approx e^{-8}$ y $\abs{u(8)}=6249$. Cerca del origen el control actúa con ganancia $x^2$ y, para vencer a la deriva $x$, hace falta $\abs u\gtrsim1/\abs x$. Por el \cref{thm:av-artstein} no existe \emph{ninguna} realimentación continua: si $\abs k\le U$ en $\Bola\delta$, para $0<x<\min\{\delta,1/U\}$ se tiene $\dot x=x(1+xk)\ge x(1-xU)>0$ y la trayectoria se aleja.
\end{ejemplo}

\codigorepo{cap10_avanzada/C10_Ej8_ControlPequeno.m}{pequeno}{$x(t)\to0$ con $\abs{u(t)}\to\infty$}

\begin{ejemplo}[label=ej:av-mano]{Diseño a mano contra Sontag}
Para $\dot x=x^3+x^2u$ (\emph{no} el ejemplo anterior), compare $u=-(1+c_1)x$ con la fórmula.
\solucion Con $V=\tfrac12x^2$, $\ell_f=x^4$, $\ell_g=x^3$. A mano: imponiendo $\Vdot=-c_1x^4$ sale $u=-(1+c_1)x$ y $\dot x=-c_1x^3$. Sontag: $k=-x(1+\sqrt{1+x^4})$ y $W=x^4\sqrt{1+x^4}$. Las dos leyes son continuas (aquí sí hay control pequeño). Cerca del origen $W\approx x^4$ y la linealización de Sontag es $-2x$: coincide con el diseño a mano con $c_1=1$. Lejos del origen Sontag usa más esfuerzo ($\abs k\ge2\abs x$) y obtiene más decrecimiento. ``Sontag requiere más esfuerzo'' solo es cierto frente al diseño con $c_1\le1$; Sontag no es la mejor ley, es una ley que no exige más ingenio una vez se tiene la CLF.
\end{ejemplo}

\subsection{Conexión con Brockett}

El integrador no holonómico $\dot x_1=u_1$, $\dot x_2=u_2$, $\dot x_3=x_1u_2-x_2u_1$ es controlable, pero no cumple la condición de \citet{Bro83}: $(x,u)\mapsto f(x,u)$ no cubre un vecindario del origen, porque para obtener $f=(0,0,\varepsilon)$ hace falta $u_1=u_2=0$, y entonces la tercera componente es $0$. No existe realimentación continua invariante en el tiempo que lo estabilice (Brockett lo probó para leyes $\Cont1$; Zabczyk y Coron lo extendieron a leyes continuas). Por \eqref{eq:av-sontagvec} y el \cref{thm:av-artstein}, \alerta{no existe una CLF suave con control pequeño}. Las salidas son realimentaciones variantes en el tiempo, discontinuas o CLF no suaves \citep[\S9.6]{KeBr23}.

\begin{pausa}
Para $\dot x=x-x^3+u$, ¿qué hace la ley de Sontag con $V=\tfrac12x^2$ cuando $\abs x$ es grande? Compárela con la ley $u=-2x+x^3$.
\respuesta{$\ell_f=x^2-x^4$, $\ell_g=x$ y $k=-x\bigl((1-x^2)+\sqrt{(1-x^2)^2+1}\bigr)$. Con $s=x^2-1$ el paréntesis vale $-s+\sqrt{s^2+1}\approx1/(2s)$, así que $k\approx-x/(2(x^2-1))\to0$ (en $x=10$, $k\approx-10/198$): Sontag conserva el término $-x^3$. La ley $u=-2x+x^3$ cancela $-x^3$, la no linealidad útil, y gasta $\approx x^3$.}
\end{pausa}

%% ---------------------------------------------------------------------
\section{El lema de Barbalat}
\label{sec:av-barbalat}
%% ---------------------------------------------------------------------

Esta sección cambia de escenario: \alerta{el lema de Barbalat es un resultado sobre funciones de una variable real (señales), no sobre ecuaciones diferenciales}. Se aplica a la señal $t\mapsto V(t,x(t))$, o a $t\mapsto W(x(t))$, sobre una solución fija. La hipótesis vive en el espacio de estados; la conclusión, en el eje del tiempo.

La situación típica es $\Vdot(t,x)\le-W(x)\le0$ con $W$ solo semidefinida positiva. En un sistema invariante en el tiempo se usaría LaSalle (capítulo~8). En uno variante en el tiempo, en general no: el conjunto $\omega$-límite de una solución de $\dot x=f(t,x)$ no tiene por qué ser invariante, porque ``las reglas del juego'' cambian con $t$. Lo que sí es gratis: si $V$ está acotada por debajo y no crece a lo largo de la solución, converge a un límite $V_\infty$, e integrando
\begin{equation}
  \int_{t_0}^\infty W\bigl(x(t)\bigr)\dif t\le V\bigl(t_0,x(t_0)\bigr)-V_\infty<\infty .
  \label{eq:av-integral}
\end{equation}

\begin{trampa}[label=tr:av-integral]{``La integral es finita, luego $W(x(t))\to0$''}
\plausible{Si el área bajo una curva no negativa es finita, la curva tiene que irse a cero.}
Una función no negativa puede tener integral finita sin tender a cero: basta con pulsos cada vez más angostos. Con pulsos triangulares centrados en $t=k$, de altura $k$ y base $1/k^3$, el área total es $\sum_k\frac1{2k^2}=\pi^2/12$, pero $\psi(k)=k\to\infty$ (\cref{fig:av-pulsos}).
\correcto{Falta una hipótesis de regularidad: la continuidad uniforme. Los pulsos tienen pendientes $\pm2k^4$, no acotadas.}
\end{trampa}

\begin{figure}[ht]
\centering
\includegraphics[width=\linewidth]{cap10/f_pulsos}
\caption{Tres trampas sobre señales. Izquierda: pulsos de área total $\pi^2/12$ con alturas que crecen sin cota. Centro: $h=\sin(\ln t)$ tiene $\dot h\to0$ pero no converge. Derecha: $h=e^{-t}\sin(e^{2t})\to0$ pero $\dot h$ no está acotada.}
\label{fig:av-pulsos}
\end{figure}

Las otras dos trampas de la figura: $\dot h\to0$ no implica que $h$ converja ($h=\sin(\ln t)$ oscila entre $\pm1$ cada vez más despacio), y $h$ convergente no implica $\dot h\to0$ ($h=e^{-t}\sin(e^{2t})$: la amplitud se apaga, pero la frecuencia crece más rápido).

\begin{definicion}[label=def:av-unifcont]{Continuidad uniforme}
$\psi:[0,\infty)\to\R$ es \clave{uniformemente continua} si para todo $\varepsilon>0$ existe $\delta(\varepsilon)>0$, \emph{el mismo para todo $t$}, tal que $\abs{t_1-t_2}<\delta\implica\abs{\psi(t_1)-\psi(t_2)}<\varepsilon$.
\end{definicion}

En la práctica casi siempre se verifica así: \alerta{si $\abs{\dot\psi(t)}\le L$ para todo $t$, entonces $\psi$ es uniformemente continua} (valor medio, $\delta=\varepsilon/L$). $\sin t$, $e^{-t}$ y $\arctan t$ lo son; $\sin(t^2)$ no (entre $t_n=\sqrt{2\pi n}$ y $s_n=\sqrt{2\pi n+\pi/2}$ la distancia tiende a cero y los valores difieren en $1$); los pulsos tampoco. Es la misma idea de ``uniforme'' de todo el capítulo: una constante que no depende del instante (\cref{def:pm-unifcont}).

\begin{lema}[label=lem:av-barbalat]{Barbălat \citep{Bar59}; \citep[Lema~8.2]{Kha02}}
\textbf{Forma integral.} Si $\psi$ es uniformemente continua y $\lim_{t\to\infty}\int_0^t\psi(\tau)\dif\tau$ existe y es finito, entonces $\psi(t)\to0$.

\textbf{Forma de derivada.} Si $h$ tiene límite finito cuando $t\to\infty$ y $\dot h$ es uniformemente continua, entonces $\dot h(t)\to0$. En particular, basta con que $\ddot h$ exista y esté acotada.
\end{lema}

\begin{demostracion}
Si $\psi\not\to0$, hay $\varepsilon>0$ y $t_k\to\infty$ con $\abs{\psi(t_k)}\ge\varepsilon$. Con $\delta=\delta(\varepsilon/2)$, en $[t_k,t_k+\delta]$ se tiene $\abs{\psi}>\varepsilon/2$ y $\psi$ no cambia de signo, así que $\bigl|\int_{t_k}^{t_k+\delta}\psi\bigr|\ge\varepsilon\delta/2$ para todo $k$. Pero si $\int_0^t\psi$ converge, el criterio de Cauchy exige que esas integrales tiendan a cero. La clave es que $\delta$ es \emph{el mismo} para todo $k$: sin continuidad uniforme, los intervalos podrían encogerse, como en los pulsos. La forma de derivada es la integral con $\psi=\dot h$.
\end{demostracion}

\begin{teorema}[label=thm:av-ly]{LaSalle--Yoshizawa \citep[Teorema~8.4]{Kha02}}
Sea $\dot x=f(t,x)$ con $f$ continua a trozos en $t$, localmente Lipschitz en $x$ uniformemente en $t$, y $f(t,0)=0$. Si $V\in\Cont1$ cumple $\alpha_1(\norm x)\le V(t,x)\le\alpha_2(\norm x)$ y $\Vdot\le-W(x)\le0$ en $[0,\infty)\times\Bola r$, con $W$ continua, entonces para $\norm{x(t_0)}<\alpha_2^{-1}(\alpha_1(r))$ las soluciones están acotadas y $W(x(t))\to0$.
\end{teorema}

\begin{demostracion}
(i) $V$ no crece, luego $x(t)$ se queda en un compacto $\Bola\rho\subset\Bola r$. (ii) Vale \eqref{eq:av-integral}. (iii) En $\Bola\rho$, $\norm{f}\le L\rho$: $x(t)$ es Lipschitz en $t$; $W$ es uniformemente continua en el compacto; la composición $t\mapsto W(x(t))$ es uniformemente continua. (iv) Barbalat, forma integral.
\end{demostracion}

\begin{lema}[label=lem:av-tipolyap]{Lema tipo Lyapunov \citep[Lema~4.3]{SlLi91}}
Sea $x(t)$ una solución y $v(t)=V(t,x(t))$. Si (1) $V$ está acotada por debajo, (2) $\dot v(t)\le0$ y (3) la \emph{señal} $t\mapsto\Vdot(t,x(t))$ es uniformemente continua, entonces $\Vdot(t,x(t))\to0$.
\end{lema}

\begin{receta}[Análisis tipo Lyapunov con Barbalat]
\begin{enumerate}
  \item Mostrar $\Vdot\le0$. Las variables en las que $V$ es definida positiva quedan acotadas.
  \item Calcular $\ddot V$ (la derivada de la señal $\dot v$) y mostrar que está acotada con las cotas del paso~1.
  \item Concluir $\Vdot\to0$ y traducirlo a las variables de estado.
\end{enumerate}
\end{receta}

\begin{trampa}[label=tr:av-barbalat]{``Barbalat da estabilidad asintótica''}
\plausible{Si $\Vdot\to0$ y $\Vdot=-2e^2$, entonces el estado converge al origen.}
El lema no pide $V$ definida positiva (basta acotada por debajo) y por eso la conclusión es \emph{parcial}: $\Vdot\to0$, no $x\to0$. No da estabilidad ni uniformidad: solo la convergencia de una señal.
\correcto{Barbalat dice qué señales convergen. Para las demás variables hace falta información adicional, como la excitación persistente del \cref{ej:av-adaptativo}. Y ``$\Vdot(t,x)$ uniformemente continua en $t$'' (como decían Handout~7 y las diapositivas) debe leerse como una propiedad de la señal $t\mapsto\Vdot(t,x(t))$, no de la función de dos variables.}
\end{trampa}

\begin{ejemplo}[label=ej:av-adaptativo]{Control adaptativo \citep[Ej.~4.13]{SlLi91}}
El error de seguimiento $e$ y el error de parámetro $\theta$ cumplen
\[
  \dot e=-e+\theta\,\omega(t),\qquad\dot\theta=-e\,\omega(t),\qquad\abs{\omega(t)}\le\bar\omega .
\]
¿Qué se puede concluir?
\solucion El sistema es variante en el tiempo (por $\omega$): LaSalle no aplica.
\emph{Paso 1.} Con $V=e^2+\theta^2$, $\Vdot=2e(-e+\theta\omega)+2\theta(-e\omega)=-2e^2\le0$. El término cruzado se cancela exactamente: \emph{para eso} se diseñó la ley $\dot\theta=-e\omega$. Así, $e$ y $\theta$ están acotadas.
\emph{Paso 2.} $\ddot V=-4e(-e+\theta\omega)$ está acotada.
\emph{Paso 3.} Por el \cref{lem:av-tipolyap}, $\Vdot=-2e^2\to0$: $e(t)\to0$.
\emph{¿Y $\theta$?} Solo se sabe que está acotada. Con $\omega=\sin t$ (excitación persistente), $\theta\to0$; con $\omega=e^{-0.3t}$, desde $(1,1)$, $\theta\to-0.11$: la regresión deja de ``ver'' el parámetro (\cref{fig:av-adaptativo}).
\end{ejemplo}

\begin{figure}[ht]
\centering
\includegraphics[width=\linewidth]{cap10/f_adaptativo}
\caption{El \cref{ej:av-adaptativo} desde $(e,\theta)=(1,1)$. En los dos casos $e\to0$, como garantiza Barbalat. Con $\omega=\sin t$, también $\theta\to0$; con $\omega=e^{-0.3t}$, $\theta$ se congela en $-0.11$.}
\label{fig:av-adaptativo}
\end{figure}

\codigorepo{cap10_avanzada/C10_Ej9_Barbalat.m}{pulsos}{Las tres trampas sobre señales}
\codigorepo{cap10_avanzada/C10_Ej10_Adaptativo.m}{adaptativo}{El error converge; el parámetro solo si hay excitación persistente}

\begin{pausa}
Con $\omega=\sin t$, ¿por qué $\theta\to0$? Use Barbalat dos veces.
\respuesta{$e\to0$ y $\ddot e=-\dot e+\dot\theta\omega+\theta\dot\omega$ está acotada (porque $\dot\omega=\cos t$ lo está), así que Barbalat aplicado a $e$ da $\dot e\to0$. De la ecuación de $e$, $\theta(t)\omega(t)=\dot e+e\to0$. Además $\dot\theta=-e\omega\to0$: $\theta$ varía cada vez más despacio. Una función que varía cada vez más despacio y cuyo producto con $\sin t$ tiende a cero no puede quedarse lejos de cero, porque en cada periodo $\sin t$ visita el valor $1$.}
\end{pausa}

%% ---------------------------------------------------------------------
\section{Estabilidad entrada--estado}
\label{sec:av-iss}
%% ---------------------------------------------------------------------

En esta sección la entrada es \emph{exógena}: no se diseña, se sufre, y se escribe $w$. Para $\dot x=Ax+Bw$ con $A$ Hurwitz y $\norm{e^{As}}\le ke^{-\lambda s}$, la variación de constantes da
\begin{equation}
  \norm{x(t)}\le\underbrace{ke^{-\lambda(t-t_0)}\norm{x(t_0)}}_{\in\KL}+\underbrace{\frac{k\norm B}{\lambda}\sup_{t_0\le\tau\le t}\norm{w(\tau)}}_{\in\K} :
  \label{eq:av-isslineal}
\end{equation}
el efecto de la condición inicial se olvida, y el de la entrada es proporcional a su tamaño. En sistemas lineales, ``GES con entrada cero'' implica ``entrada acotada, estado acotado''. En no lineales, no.

\begin{ejemplo}[label=ej:av-riccati]{GES con $w=0$ y escape con $w=1$}
Estudie $\dot x=-3x+(1+2x^2)w$ con $w\equiv1$ y $x(0)=2$.
\solucion Con $w=0$ el origen es GES. Con $w\equiv1$ la ecuación es de Riccati y factoriza: $\dot x=2x^2-3x+1=(2x-1)(x-1)$. Por fracciones parciales, $\frac{x-1}{2x-1}=\frac{e^t}{3}$ y
\[
  x(t)=\frac{3-e^t}{3-2e^t} ,
\]
cuyo denominador se anula en $t_e=\ln\tfrac32=0.405$: \alerta{escape en tiempo finito con una entrada constante de tamaño $1$}. (Handout~7 tenía $3-e^{-t}$ en el numerador.) Más simple aún: $\dot x=-x+xw$ con $w\equiv c>1$ crece exponencialmente.
\end{ejemplo}

\begin{definicion}[label=def:av-iss]{ISS \citep{Son89a}; \citep[Def.~4.7]{Kha02}}
$\dot x=f(t,x,w)$, con $f$ continua a trozos en $t$ y localmente Lipschitz en $(x,w)$, es \clave{estable entrada--estado (ISS)} si existen $\beta\in\KL$ y $\gamma\in\K$ tales que, para toda condición inicial y toda entrada $w$ continua a trozos y acotada, la solución existe para todo $t\ge t_0$ y
\begin{equation}
  \norm{x(t)}\le\beta\bigl(\norm{x(t_0)},t-t_0\bigr)+\gamma\Bigl(\sup_{t_0\le\tau\le t}\norm{w(\tau)}\Bigr)\qquad\forall\,t\ge t_0\ge0 .
  \label{eq:av-iss}
\end{equation}
\end{definicion}

$\beta$ es el \emph{transitorio}, la memoria de la condición inicial; $\gamma$ es la \clave{ganancia ISS}, cuánto estado compra una entrada de tamaño dado, y no tiene por qué ser lineal. La versión con $\max\{\hat\beta,\hat\gamma\}$ es equivalente, porque $a+b\le\max\{2a,2b\}\le2a+2b$. El supremo es \emph{causal}. Y la definición pide que la solución exista para todo $t$: el \cref{ej:av-riccati} muestra que eso no es gratis.

\begin{proposicion}[label=prop:av-issconsec]{Tres consecuencias de ISS}
Si $\dot x=f(t,x,w)$ es ISS: (1) con $w\equiv0$ el origen es GUAS; (2) si $w$ es acotada, $\limsup\norm{x(t)}\le\gamma(\norm w_\infty)$ y las soluciones son GUUB con cualquier cota $b>\gamma(\norm w_\infty)$; (3) si $w(t)\to0$, entonces $x(t)\to0$.
\end{proposicion}

\begin{demostracion}
(1) y (2) se leen en \eqref{eq:av-iss}. (3) \emph{Reiniciar el reloj:} dado $\varepsilon$, sea $T_1$ con $\gamma(\norm{w(t)})\le\varepsilon/2$ para $t\ge T_1$. Como la definición vale para todo $t_0$, desde $T_1$, $\norm{x(t)}\le\beta(\norm{x(T_1)},t-T_1)+\varepsilon/2$, y basta esperar a que el primer término baje de $\varepsilon/2$ \citep[Lec.~14]{Arc18}.
\end{demostracion}

\begin{teorema}[label=thm:av-isslyap]{Funciones ISS-Lyapunov \citep[Teorema~4.19]{Kha02}}
Sea $V:[0,\infty)\times\R^n\to\R$ de clase $\Cont1$ con
\[
  \alpha_1(\norm x)\le V(t,x)\le\alpha_2(\norm x),\qquad\norm x\ge\rho(\norm w)\implica\jac Vt+\jac Vxf(t,x,w)\le-W_3(x),
\]
$\alpha_1,\alpha_2\in\Kinf$, $\rho\in\K$ y $W_3$ continua y definida positiva. Entonces el sistema es ISS con
\begin{equation}
  \gamma=\alpha_1^{-1}\circ\alpha_2\circ\rho .
  \label{eq:av-gananciaiss}
\end{equation}
\end{teorema}

\begin{demostracion}
\alerta{Es el teorema de UUB disfrazado.} Fije $w$ y $t$, y sea $\bar w=\sup_{[t_0,t]}\norm w$. Para esa entrada, $\dot x=f(s,x,w(s))$ es un sistema variante en el tiempo \emph{sin entrada}, con $\Vdot\le-W_3$ para $\norm x\ge\mu=\rho(\bar w)$: es la hipótesis del \cref{thm:av-uub} en forma global. Luego $\norm{x(s)}\le\max\{\beta(\norm{x(t_0)},s-t_0),\ \alpha_1^{-1}(\alpha_2(\rho(\bar w)))\}$. La $\beta$ del \cref{thm:av-uub} se construyó con $\alpha_1,\alpha_2,\alpha_3$ y no depende de $\mu$: es la misma para todas las entradas.
\end{demostracion}

ISS es UUB con un radio que depende de la entrada, y la \cref{tr:av-cota} reaparece: la ganancia no es $\rho$ sino $\alpha_1^{-1}\circ\alpha_2\circ\rho\ge\rho$. (Handout~7 pedía $\rho\in\Kinf$ y $W_3\in\K$; bastan $\rho\in\K$ y $W_3$ definida positiva.)

\begin{teorema}[label=thm:av-sontagwang]{Sontag--Wang \citep{SoWa95}; \citep[Teorema~7.4]{KeBr23}}
$\dot x=f(x,w)$ es ISS si y solo si existen $V\in\Cont1$ y $\alpha_1,\alpha_2,\alpha_3,\chi\in\Kinf$ con $\alpha_1(\norm x)\le V(x)\le\alpha_2(\norm x)$ y $\norm x\ge\chi(\norm w)\implica\nabla V\cdot f(x,w)\le-\alpha_3(\norm x)$. Equivalentemente (\clave{forma disipativa}), existen $\alpha_3,\sigma\in\Kinf$ con
\begin{equation}
  \nabla V(x)\cdot f(x,w)\le-\alpha_3(\norm x)+\sigma(\norm w) .
  \label{eq:av-issdisip}
\end{equation}
\end{teorema}

En la forma disipativa, $V$ es un almacenamiento que se disipa a ritmo $\alpha_3$ y en el que la entrada solo puede inyectar $\sigma(\norm w)$: es la antesala del \cref{cap:pasividad}. De la forma disipativa a la de implicación hay un paso: si $\norm x\ge\alpha_3^{-1}(2\sigma(\norm w))$, $\Vdot\le-\tfrac12\alpha_3(\norm x)$. Lo difícil es ``ISS $\implica$ existe $V$'', que es un converso.

\begin{ejemplo}[label=ej:av-issej]{Dos funciones ISS-Lyapunov}
Muestre que (a) $\dot x=-x^3+w$ y (b) $\dot x=-x-2x^3+(1+x^2)w^2$ son ISS y calcule sus ganancias.
\solucion (a) Con $V=\tfrac12x^2$ y $0<\theta<1$, $\Vdot=-(1-\theta)x^4-\theta x^4+xw\le-(1-\theta)x^4$ si $\abs x\ge(\abs w/\theta)^{1/3}$. Como $\alpha_1=\alpha_2$, $\gamma(r)=(r/\theta)^{1/3}$; con $\theta=\tfrac12$, $\gamma(0.5)=1$. La ganancia no es lineal: entradas pequeñas producen estados relativamente grandes ($r^{1/3}\gg r$ si $r\ll1$), el precio de no tener parte lineal en la disipación. Con $w=0.5\sin2t$, $\limsup\abs x$ queda por debajo de $1$.
(b) Si $\abs x\ge w^2$, $x(1+x^2)w^2\le\abs x(1+x^2)\abs x=x^2+x^4$ y $\Vdot\le-x^4$: $\rho(r)=r^2=\gamma(r)$. (Handout~7 partía de ``$u^2\le x$''; debe ser $w^2\le\abs x$, porque si $x<0$ la otra desigualdad es imposible.)
\end{ejemplo}

\codigorepo{cap10_avanzada/C10_Ej11_ISScontraejemplos.m}{contra}{GES con entrada cero y escape con entrada acotada}
\codigorepo{cap10_avanzada/C10_Ej12_ISSenvolvente.m}{envolvente}{La envolvente $\beta+\gamma$ y el $\limsup$ simulado}

Un atajo útil \citep[Lema~4.6]{Kha02}: si $f(t,x,w)$ es $\Cont1$ y \emph{globalmente} Lipschitz en $(x,w)$ y el origen de $\dot x=f(t,x,0)$ es GES, el sistema es ISS. La prueba es la del \cref{lem:av-robustez}: el converso da $\norm{\partial V/\partial x}\le c_4\norm x$ y $\Vdot\le-c_3\norm x^2+c_4L\norm x\norm w$. Sin Lipschitz global no hay atajo: el término $2x^2w$ del \cref{ej:av-riccati} lo rompe. Para la familia $\dot x=-x^r+x^sw$, con $r$ impar y $s\ge0$ enteros, hay una regla limpia: \alerta{es ISS si y solo si $r>s$} \citep[Lec.~14]{Arc18}; la disipación debe crecer más rápido que el acoplamiento con la entrada. (Si $s\ge r$, con $w\equiv c>1$ y $x(0)$ grande, $\dot x=x^r(cx^{s-r}-1)>0$.)

\subsection{Estabilizar no es robustecer}

\begin{ejemplo}[label=ej:av-rediseno]{Cancelación exacta frente a amortiguamiento \citep[Ej.~9.2]{KeBr23}}
Para $\dot x=x+(x^2+1)u$, la ley $k(x)=-2x/(x^2+1)$ da $\dot x=-x$: GES. ¿Qué pasa con un error aditivo en el actuador, $u=k(x)+w$?
\solucion El lazo queda $\dot x=-x+(x^2+1)w$, que no es ISS. Con $w\equiv1$ y $x(0)=0$, $\dot x=x^2-x+1\ge\tfrac34$ y
\[
  t_e=\int_0^\infty\frac{\dif x}{x^2-x+1}=\frac{4\pi}{3\sqrt3}=2.42 :
\]
\alerta{se escapa incluso partiendo del equilibrio}. Con $w\equiv\varepsilon\in(0,\tfrac12]$, $\varepsilon x^2-x+\varepsilon=0$ tiene raíces $x_\pm=(1\pm\sqrt{1-4\varepsilon^2})/(2\varepsilon)$; con $\varepsilon=0.1$, $x_-=0.101$ y $x_+=9.90$. Desde $x(0)=0$ la solución converge a $x_-$, pero desde cualquier $x(0)>x_+\approx1/\varepsilon$ escapa: \emph{por pequeña que sea la perturbación, hay condiciones iniciales que escapan}.
\emph{Rediseño.} Con $V_1=\tfrac12x^2$, $\Lie g{V_1}=x(x^2+1)$ y $k_1=k-\Lie g{V_1}$:
\[
  \dot x=-x-x(x^2+1)^2+(x^2+1)w,\qquad\Vdot_1\le-x^2-\tfrac12x^2(x^2+1)^2+\tfrac12w^2 ,
\]
por Young, que es la forma disipativa \eqref{eq:av-issdisip}: el lazo rediseñado es ISS. Con $w\equiv1$ y $x(0)=0$, $x(t)$ se asienta en el equilibrio $0.4887$ (\cref{fig:av-rediseno}).
\end{ejemplo}

\begin{figure}[ht]
\centering
\includegraphics[width=\linewidth]{cap10/f_rediseno}
\caption{El \cref{ej:av-rediseno} con $w\equiv1$ desde $x(0)=0$. La ley de cancelación exacta (coral) escapa en $t=2.42$; la ley con amortiguamiento $-\Lie gV$ (azul) se asienta en $x=0.4887$. A la derecha, el lazo cancelado con $w\equiv0.1$: $x_-=0.101$ atrae, y las condiciones iniciales por encima de $x_+=9.90$ escapan.}
\label{fig:av-rediseno}
\end{figure}

\codigorepo{cap10_avanzada/C10_Ej13_RedisenoISS.m}{rediseno}{Cancelación exacta contra amortiguamiento $-\Lie gV$}

\video{V10_4_rediseno_iss}{Estabilizar no es robustecer}

El procedimiento general \citep[\S9.2]{KeBr23}: si $k$ estabiliza con función de Lyapunov $V$ (que existe por un converso), entonces $u=k(x)-\Lie gV(x)$ es ISS respecto a errores aditivos en el actuador, porque
\[
  \Vdot=\Lie fV+\Lie gV\,k-(\Lie gV)^2+\Lie gV\,w\le-\alpha(\norm x)-\tfrac12(\Lie gV)^2+\tfrac12w^2 .
\]
Es el \emph{amortiguamiento no lineal} (control de Jurdjevic--Quinn~\citep{JuQu78}), el mismo término que el $\kappa$ de \eqref{eq:av-sontagvec}. Con Sontag y $\kappa>0$, $\Vdot\le-W(x)-\frac\kappa2\norm{\ell_g}^2+\frac1{2\kappa}\norm w^2$, que es ISS \emph{siempre que} $W\ge\alpha_3(\norm x)$ con $\alpha_3\in\Kinf$; esa condición hay que verificarla en cada caso.

\begin{trampa}[label=tr:av-robustecer]{``Si el lazo es GES, aguanta perturbaciones pequeñas''}
\plausible{Un lazo cerrado exponencialmente estable tiene margen: una perturbación pequeña solo lo desplaza un poco.}
Eso es cierto localmente (y globalmente si todo es globalmente Lipschitz). En el \cref{ej:av-rediseno} el lazo es GES y lineal, pero la perturbación entra multiplicada por $x^2+1$: lejos del origen ``pequeña'' deja de serlo.
\correcto{La cancelación exacta deja el lazo frágil. Robustecer es una exigencia de diseño aparte, y se logra dominando el término cruzado con $-\Lie gV$.}
\end{trampa}

\begin{pausa}
Muestre que $V_2=\tfrac12\ln(1+x^2)$ también es una función ISS-Lyapunov para $\dot x=-x-x(x^2+1)+(x^2+1)w$. ¿Qué ventaja tiene sobre $V_1$?
\respuesta{$\nabla V_2=x/(1+x^2)$, así que $\Lie g{V_2}=x$ y $\Vdot_2=-\frac{x^2}{1+x^2}-x^2+xw\le-\frac{x^2}{1+x^2}-\frac12x^2+\frac12w^2$. $V_2$ es de clase $\Kinf$ en $\abs x$, luego es ISS-Lyapunov. Con $V_2$ el amortiguamiento que hace falta es $-x$, no $-x(x^2+1)$: la escala de $V$ cambia cuánto control se gasta, como en la \cref{tr:av-escala}.}
\end{pausa}

%% ---------------------------------------------------------------------
\section{Sistemas interconectados}
\label{sec:av-inter}
%% ---------------------------------------------------------------------

En esta sección $x_1\in\R^{n_1}$ y $x_2\in\R^{n_2}$ son los estados de \emph{dos subsistemas}; en los ejemplos son escalares, pero el razonamiento es por subsistemas. Para la cascada
\[
  \dot x_1=f_1(x_1,x_2),\qquad\dot x_2=f_2(x_2),
\]
con el origen de $\dot x_2=f_2(x_2)$ GAS y el de $\dot x_1=f_1(x_1,0)$ GAS, la intuición dice: ``como $x_2\to0$, a la larga $x_1$ ve su sistema sin entrada''. El problema está en ``a la larga'': antes de que $x_2$ se apague, $x_1$ puede escapar en tiempo finito, y entonces no hay ``larga''.

\begin{ejemplo}[label=ej:av-cascada]{GAS más GAS no da GAS \citep[Ej.~4]{Arc18}}
Estudie $\dot x_1=-x_1+x_1^2x_2$, $\dot x_2=-x_2$.
\solucion Con $x_2=x_{20}e^{-t}$ y $z=1/x_1$, $\dot z=z-x_{20}e^{-t}$, lineal, con solución $z=(z_0-x_{20}/2)e^t+(x_{20}/2)e^{-t}$. $x_1$ escapa cuando $z$ se anula, lo que ocurre en
\[
  t_e=\frac12\ln\frac{x_{10}x_{20}}{x_{10}x_{20}-2}
\]
si y solo si $x_{10}x_{20}>2$. Desde $(2,2)$, $t_e=\tfrac12\ln2=0.347$ (\cref{fig:av-cascada}).
\end{ejemplo}

\begin{trampa}[label=tr:av-cascadaerr]{``$\dot x_1=-x_1+x_1x_2^2$ también escapa''}
\plausible{Si $x_1^2x_2$ produce escape, $x_1x_2^2$ debería producirlo también; así lo decían Handout~7 y las diapositivas.}
Es falso. La ecuación de $x_1$ es lineal en $x_1$ y
\[
  x_1(t)=x_{10}\exp\Bigl(-t+\frac{x_{20}^2}{2}\bigl(1-e^{-2t}\bigr)\Bigr)\to0 ,
\]
con $\abs{x_1(t)}\le\abs{x_{10}}e^{x_{20}^2/2}$: el origen es GAS. El transitorio puede ser grande (desde $(1,3)$, $x_1$ sube hasta $18.2$), pero no escapa.
\correcto{Sin embargo el subsistema $x_1$ no es ISS respecto a $x_2$: con $x_2\equiv c>1$, $\dot x_1=(c^2-1)x_1$ crece. ISS de $x_1$ es \emph{suficiente} para que la cascada sea GAS, no necesaria.}
\end{trampa}

\begin{proposicion}[label=prop:av-cascadaiss]{Cascadas ISS \citep{Son89a}; \citep[Teorema~7.9]{KeBr23}}
Si en $\dot x_1=f_1(x_1,x_2)$, $\dot x_2=f_2(x_2,w)$ el subsistema $x_1$ es ISS con $x_2$ como entrada y el subsistema $x_2$ es ISS con $w$ como entrada, la cascada es ISS con entrada $w$.
\end{proposicion}

\begin{corolario}[label=cor:av-cascadagas]{\citep[Lema~4.7]{Kha02}}
Si el origen de $\dot x_2=f_2(x_2)$ es GUAS y $\dot x_1=f_1(x_1,x_2)$ es ISS respecto a $x_2$, el origen de la cascada es GUAS.
\end{corolario}

\begin{demostracion}
Por ISS de $x_1$, para todo $s\in[t_0,t]$,
$\norm{x_1(t)}\le\beta_1(\norm{x_1(s)},t-s)+\gamma_1\bigl(\sup_{s\le\tau\le t}\norm{x_2(\tau)}\bigr)$, y $\norm{x_2(\tau)}\le\beta_2(\norm{x_2(t_0)},\tau-t_0)$. Con $s=t_0$ sale el acotamiento (y la estabilidad): $\norm{x_1(t)}\le\beta_1(\norm{x_1(t_0)},0)+\gamma_1(\beta_2(\norm{x_2(t_0)},0))$. Con $s=(t+t_0)/2$, \emph{el punto medio},
\[
  \norm{x_1(t)}\le\beta_1\Bigl(\norm{x_1(s)},\tfrac{t-t_0}2\Bigr)+\gamma_1\Bigl(\beta_2\bigl(\norm{x_2(t_0)},\tfrac{t-t_0}2\bigr)\Bigr)\to0 :
\]
la mitad del tiempo para que $x_2$ se apague y la otra mitad para que $x_1$ olvide su condición en $s$.
\end{demostracion}

\begin{figure}[ht]
\centering
\includegraphics[width=\linewidth]{cap10/f_cascada}
\caption{Izquierda: la cascada del \cref{ej:av-cascada} desde $(2,2)$ escapa en $t=0.347$. Centro: $\dot x_1=-x_1+x_1x_2^2$ desde $(1,3)$ sube hasta $18.2$ y vuelve a cero. Derecha: la cascada ISS del \cref{ej:av-cascadaiss} con $w=0.3\sin t$, acotada.}
\label{fig:av-cascada}
\end{figure}

\begin{ejemplo}[label=ej:av-cascadaiss]{$\dot x_1=-x_1+x_2^2$, $\dot x_2=-x_2+w$}
Pruebe que es ISS de dos maneras: por cascada y con una función de Lyapunov compuesta.
\solucion \emph{Por cascada.} Con $V=\tfrac12x_1^2$, $\Vdot=-x_1^2+x_1x_2^2\le-(1-\theta)x_1^2$ si $\abs{x_1}\ge x_2^2/\theta$: el subsistema $x_1$ es ISS respecto a $x_2$ con $\gamma_1(r)=r^2/\theta$; el $x_2$ lo es respecto a $w$. Por la \cref{prop:av-cascadaiss}, la cascada es ISS.
\emph{Con $V$ compuesta.} Se prueba $V=\tfrac12x_1^2+\tfrac a4x_2^4$ (cuártica en $x_2$ porque el acoplamiento $x_1x_2^2$ es de orden~4 al completar el cuadrado). Con $w=0$,
\[
  \Vdot=-\Bigl(x_1-\frac{x_2^2}2\Bigr)^2-\Bigl(a-\frac14\Bigr)x_2^4 ,
\]
negativa definida si $a>\tfrac14$. Con $a=2$ y $w$, $\Vdot=-x_1^2+x_1x_2^2-2x_2^4+2x_2^3w$, y dos usos de Young \eqref{eq:av-young}: $x_1x_2^2\le\tfrac12x_1^2+\tfrac12x_2^4$ y, con $p=\tfrac43$, $q=4$, $\lambda^4=8/27$, $\abs{x_2}^3\abs w\le\tfrac12x_2^4+\tfrac{27}{32}w^4$. Queda
\[
  \Vdot\le-\tfrac12x_1^2-\tfrac12x_2^4+\tfrac{27}{16}w^4 ,
\]
la forma disipativa: ISS. (Con $a=1$, Handout~7 tenía un signo duplicado; lo correcto es $\Vdot=-\tfrac12(x_1-x_2^2)^2-\tfrac12(x_1^2+x_2^4)+x_2^3w$.)
\end{ejemplo}

\codigorepo{cap10_avanzada/C10_Ej14_Cascadas.m}{cascadas}{Escape, transitorio grande y cascada ISS}

\subsection{Retroalimentación y pequeña ganancia}

Si en lugar de una cascada hay una \emph{retroalimentación}, $\dot x_1=f_1(x_1,x_2)$, $\dot x_2=f_2(x_2,x_1)$, con cada subsistema ISS respecto al otro (ganancias $\gamma_1$, $\gamma_2$), la ISS de cada uno ya no basta: la señal da vueltas por el lazo y puede amplificarse.

\begin{teorema}[label=thm:av-pequenaganancia]{Pequeña ganancia \citep{JTP94}}
Si $\gamma_1\circ\gamma_2(r)<r$ para todo $r>0$, el origen de la interconexión es GAS.
\end{teorema}

La parte fácil de la prueba muestra la idea. Si las soluciones están acotadas y $L_i=\limsup\norm{x_i(t)}$, reiniciando el reloj como en la \cref{prop:av-issconsec}, $L_1\le\gamma_1(L_2)$ y $L_2\le\gamma_2(L_1)$, así que $L_1\le\gamma_1(\gamma_2(L_1))<L_1$ si $L_1>0$: imposible. \alerta{Una vuelta completa por el lazo contrae.} Lo difícil es probar el acotamiento y la estabilidad.

\begin{ejemplo}[label=ej:av-pequenaganancia]{$\dot x_1=-x_1+a\,x_2$, $\dot x_2=-x_2+b\sin x_1$}
¿Para qué $a,b$ es GAS el origen?
\solucion Por variación de constantes, $\gamma_1(r)=\abs ar$, y como $\abs{\sin x_1}\le\abs{x_1}$, $\gamma_2(r)=\abs br$. La condición es $\abs{ab}<1$: GAS.
\emph{Con $ab>1$.} Los equilibrios cumplen $x_1=ab\sin x_1$, con soluciones $\pm\eq_1\neq0$; con $a=1$, $b=1.2$, $\eq_1=1.027$. El jacobiano tiene valores propios $-1\pm\sqrt{ab\cos x_1}$: en el origen, $-1\pm\sqrt{1.2}$, silla; en $\pm\eq_1$, $ab\cos\eq_1=0.62<1$, nodos estables. El origen deja de ser GAS.
\emph{Con $ab=-3$: GAS aunque $\abs{ab}>1$.} $x_1=-3\sin x_1$ solo tiene la solución $0$; las soluciones son últimamente acotadas; la divergencia es $-2$, así que no hay órbitas cerradas (Bendixson, \cref{cap:orbitas}); el único equilibrio es un foco estable ($-1\pm j\sqrt3$), y por Poincaré--Bendixson es el $\omega$-límite de todas las trayectorias. \alerta{La pequeña ganancia es suficiente, no necesaria.}
\end{ejemplo}

\begin{figure}[ht]
\centering
\includegraphics[width=\linewidth]{cap10/f_pequena_ganancia}
\caption{El \cref{ej:av-pequenaganancia}. Izquierda: $ab=0.8$, todas las trayectorias van al origen. Centro: $ab=1.2$, el origen es una silla y aparecen dos nodos estables en $x_1=\pm1.027$. Derecha: $ab=-3$, viola la pequeña ganancia y aun así el origen es GAS (foco).}
\label{fig:av-pequenaganancia}
\end{figure}

\codigorepo{cap10_avanzada/C10_Ej15_PequenaGanancia.m}{ganancia}{Los tres retratos de fase de la retroalimentación}

\begin{pausa}
$\dot x_1=-x_1+2x_2$, $\dot x_2=-x_2+x_1$: ¿cada subsistema es ISS respecto al otro? ¿Se cumple la pequeña ganancia? ¿Es GAS el origen?
\respuesta{Cada subsistema es lineal estable y, por \eqref{eq:av-isslineal}, ISS con ganancias $2$ y $1$. Pero $\gamma_1\gamma_2=2>1$: la pequeña ganancia falla. Y el origen no es GAS: la matriz $\bigl[\begin{smallmatrix}-1&2\\1&-1\end{smallmatrix}\bigr]$ tiene el valor propio $-1+\sqrt2>0$. Aquí la falla de la condición sí se traduce en inestabilidad.}
\end{pausa}

%% ---------------------------------------------------------------------
\section{Para llevar}
\label{sec:av-cierre}
%% ---------------------------------------------------------------------

\begin{parallevar}
\begin{enumerate}
  \item \textbf{UUB:} sin equilibrio se garantiza ``llegar a una bola y quedarse'', con constantes independientes de $t_0$. La herramienta es el argumento de los dos conjuntos de nivel, y la cota última es $\alpha_1^{-1}(\alpha_2(\mu))$, no $\mu$.
  \item \textbf{Conversos:} la estabilidad implica que $V$ existe, pero no la construye. Su valor está en las propiedades cuantitativas ($c_4$, la cota uniforme $\alpha_3$) que permiten probar robustez. ``Solo estable'' puede exigir una $V(t,x)$ aunque el sistema sea invariante en el tiempo.
  \item \textbf{CLF y Sontag:} con una CLF, una fórmula universal da la ley de control; la continuidad en el origen equivale a la propiedad de control pequeño (Artstein). La fórmula no gasta control donde la deriva ayuda, y la escala de $V$ es una perilla.
  \item \textbf{Barbalat:} $\Vdot\le0$ y $\ddot V$ acotada implican $\Vdot\to0$. La conclusión es parcial y vive en el eje del tiempo; sirve donde LaSalle no aplica, en sistemas variantes en el tiempo.
  \item \textbf{ISS:} es UUB con un radio que depende de la entrada, con ganancia $\alpha_1^{-1}\circ\alpha_2\circ\rho$. Estabilizar no es robustecer: la cancelación exacta se arregla con amortiguamiento $-\Lie gV$.
  \item \textbf{Interconexiones:} las cascadas ISS se componen; GAS más GAS no basta. Las retroalimentaciones necesitan pequeña ganancia, que es suficiente pero no necesaria.
\end{enumerate}
\end{parallevar}

\begin{trampa}[title={\faBug\hspace{0.45em}Las trampas del capítulo, en una caja}]{}
\begin{itemize}
  \item UUB no es estabilidad con otra cota (\cref{tr:av-noestable}); la cota última no es $\mu$ (\cref{tr:av-cota}); el $T$ de Lyapunov es una cota pesimista (\cref{tr:av-T}).
  \item El converso no es constructivo (\cref{tr:av-converso}).
  \item Una $V$ puede no ser CLF aunque el sistema sea estabilizable (\cref{tr:av-noclf}); el exponente 4 no es decorativo (\cref{tr:av-cuarta}); ``universal'' no es ``sin perillas'' (\cref{tr:av-escala}); si $f(0)\neq0$, el control en régimen no es cero (\cref{tr:av-regimen}).
  \item $\int W<\infty$ no implica $W\to0$ (\cref{tr:av-integral}); Barbalat no da estabilidad asintótica (\cref{tr:av-barbalat}).
  \item ``Existe $u$'' (CLF) no es ``para todo $w$'' (ISS); un lazo GES puede no aguantar perturbaciones pequeñas (\cref{tr:av-robustecer}).
  \item GAS más GAS no da GAS en cascada (\cref{ej:av-cascada}), pero no toda cascada no ISS escapa (\cref{tr:av-cascadaerr}).
\end{itemize}
\end{trampa}

\begin{profundiza}[Correcciones respecto a la versión anterior (Handout7 y las diapositivas de estabilidad avanzada)]
(i) La cota última del \cref{thm:av-uub} es $\alpha_1^{-1}(\alpha_2(\mu))$ y la igualdad con $\mu$ exige $\alpha_1(\mu)=\alpha_2(\mu)$, no $\alpha_1=\alpha_2$; el comportamiento $\KL$ vale hasta que la trayectoria entra en $\Nivel\eta$, no ``mientras $\norm x\ge\mu$''. (ii) En los corolarios UB/UUB aparecía $\Bola{\alpha^{-1}(\eta)}$ por $\Bola{\alpha_1^{-1}(\eta)}$, faltaba el conjunto de condiciones iniciales y la hipótesis se escribía ``$\dot V(x(t))\le-W_3(x)$'' en lugar de $\partial V/\partial t+(\partial V/\partial x)f\le-W_3(x)$. (iii) El converso para UAS da $\Vdot\le-\alpha_3(\norm x)$, una cota uniforme, no solo $\Vdot<0$. (iv) Teorema de Artstein: el nombre es Artstein (no ``Arstein''), y faltaban la no acotación radial y la regularidad de $k$. (v) Fórmula de Sontag: $W(x)=\sqrt{\ell_f^2+\ell_g^4}$, no ``$W(x)\sqrt{\cdots}$''; el término $-\kappa(\Lie gV)^2$ está en $\Vdot$, mientras en la ley de control el término es $-\kappa\Lie gV$; la versión vectorial \eqref{eq:av-sontagvec} hace falta para Brockett, donde $u\in\R^2$, y la extensión a leyes continuas es de Zabczyk y Coron. (vi) La gráfica de $u(t)$ del ejemplo $\cos x-x^3+u$ arrancaba en $\approx-20$ y tendía a $0$; con la ley correcta arranca en $0.17$ y tiende a $-1$. (vii) La diapositiva del diseño a mano no decía el sistema: es $\dot x=x^3+x^2u$, no $\dot x=x+x^2u$; ``Sontag requiere más esfuerzo'' solo vale frente a $c_1\le1$. (viii) Lema tipo Lyapunov: la continuidad uniforme es de la señal $t\mapsto\Vdot(t,x(t))$. En Kellett--Braun, Barbalat es el Lema~A.12 (el~A.4 es la desigualdad de Young). (ix) ISS lineal: $\lambda>0$, la integral relevante es $\int e^{-\lambda(t-\tau)}\dif\tau$ y faltaba $\norm B$. El contraejemplo de Riccati tiene numerador $3-e^t$, no $3-e^{-t}$. (x) Teorema ISS-Lyapunov: bastan $\rho\in\K$ y $W_3$ definida positiva, y las cotas de $V$ se llaman $\alpha_1,\alpha_2$ (no $\gamma_1,\gamma_2$, que chocan con la ganancia). En el ejemplo $\dot x=-x-2x^3+(1+x^2)w^2$ la condición es $w^2\le\abs x$, no $u^2\le x$. (xi) La regla de Arcak solo da la dirección ``si''; la recíproca se agregó. (xii) La cascada $\dot x_1=-x_1+x_1x_2^2$ \emph{no} escapa (\cref{tr:av-cascadaerr}). (xiii) La proposición ISS más ISS es de Sontag (Teorema~7.9 de Kellett--Braun); el Lema~4.7 de Khalil es el corolario GUAS más ISS. En la $V$ compuesta con $a=1$ había un signo duplicado. (xiv) La predicción~3 se respondía con $w\equiv1$, que no es pequeña; ahora se responde con $w\equiv\varepsilon$ y el umbral $x_+\approx1/\varepsilon$. (xv) Notación: entrada exógena $w$ y control $u$; $\ell_f,\ell_g$ en lugar de $a,b$ (que chocaban con el radio inicial y la cota última); $k,\lambda$ y la ventana $h$ en el converso; $\nu$ para la cota de la perturbación; $h,\psi$ para señales; ``invariante en el tiempo'' en lugar de ``autónomo''. (xvi) Los conversos, que la planeación anterior asignaba al capítulo~9, quedaron en este capítulo.
\end{profundiza}

%% ---------------------------------------------------------------------
\section{Problemas}
\label{sec:av-problemas}
%% ---------------------------------------------------------------------

Los problemas marcados con \etiquetaM{} requieren \matlab/\simulink; los marcados con \etiquetaEstrella{} son de profundización. Cada uno tiene una trampa: búsquela antes de dar la respuesta.

\subsection*{Acotamiento último}

\begin{problema}[cuándo la cota es $\mu$]\label{prob:av-cubica}
Para $\dot x=-x^3+\delta\sin t$: (a) con $V=\tfrac12x^2$, encuentre $\mu$ y la cota última garantizada; (b) ¿por qué aquí $b=\mu$ y en el \cref{ej:av-oscilador} no? (c) \etiquetaM{} con $\delta=0.5$, compare la cota con el $\limsup$ simulado.
\end{problema}

\begin{problema}[cuánto se espera]\label{prob:av-esperaresorte}
Para el resorte del \cref{ej:av-resorte} con $M=0.3$, $\theta=0.9$ y $x(t_0)=(2,0)$: (a) calcule $c=V(x(t_0))$, $\eta=\alpha_2(\mu)$ y el ritmo mínimo $k$ en la corona; (b) acote el tiempo de espera $T$; (c) \etiquetaM{} compare con el tiempo que tarda la simulación en entrar a $\Nivel\eta$.
\end{problema}

\begin{problema}[acotado sin cota última]\label{prob:av-integrador}
Para $\dot x=\delta\sin t$: (a) muestre que las soluciones son UB; (b) muestre que no son UUB; (c) ¿qué hipótesis del \cref{thm:av-uub} falla y por qué no se puede arreglar cambiando $V$?
\end{problema}

\subsection*{Conversos y CLF}

\begin{problema}[un margen garantizado]\label{prob:av-margen}
Para $A=\bigl[\begin{smallmatrix}0&1\\-2&-3\end{smallmatrix}\bigr]$: (a) resuelva $A\T P+PA=-I$; (b) calcule el margen $\nu^*$ del \cref{lem:av-robustez}; (c) encuentre una perturbación lineal $g_p$ con $\norm{g_p}\le\nu\norm x$ que desestabilice el sistema. ¿Para qué $\nu$ lo logra? ¿Es mayor que $\nu^*$?
\end{problema}

\begin{problema}[arreglar la $V$]\label{prob:av-dobleintegrador}
Para el doble integrador $\dot x_1=x_2$, $\dot x_2=u$: (a) muestre que $V=x_1^2+x_1x_2+x_2^2$ es CLF; (b) escriba la ley de Sontag y estudie $k(sx)/s$ cuando $s\to0^+$; (c) ¿tiene el lazo cerrado una linealización en el origen?
\end{problema}

\begin{problema}[conservar lo útil]\label{prob:av-conservar}
Para $\dot x=x-x^3+u$: (a) escriba la ley de Sontag con $V=\tfrac12x^2$; (b) \etiquetaM{} desde $x(0)=3$, compare el esfuerzo $\int_0^{10}u^2\dif t$ con el de $u=-2x+x^3$; (c) ¿cuál converge más rápido? ¿Cómo puede la ley que gasta menos converger más rápido?
\end{problema}

\subsection*{Barbalat}

\begin{problema}[control adaptativo por modelo de referencia]\label{prob:av-mrac}
Planta $\dot y=a^*y+u$ con $a^*$ desconocido; modelo $\dot y_m=-ay_m+r(t)$, $a>0$, $r$ acotada; ley $u=-k(t)y+r$ y adaptación $\dot k=ey$, con $e=y-y_m$. (a) Con $k^*=a^*+a$ y $\tilde k=k-k^*$, muestre que $\dot e=-ae-\tilde ky$; (b) con $V=\tfrac12e^2+\tfrac12\tilde k^2$, pruebe que $e\to0$; (c) ¿converge $k$ a $k^*$?
\end{problema}

\begin{problema}[cuando falta la continuidad uniforme]\label{prob:av-sint2}
Sea $h(t)=\sin(t^2)/t$. (a) Muestre que $h\to0$; (b) muestre que $\dot h\not\to0$; (c) ¿qué hipótesis de la forma de derivada del \cref{lem:av-barbalat} falla?
\end{problema}

\subsection*{ISS e interconexiones}

\begin{problema}[ISS o no]\label{prob:av-issono}
Decida si son ISS y, cuando lo sean, dé una ganancia: (a) $\dot x=-x+x^2w$; (b) $\dot x=-x^3+x^2w$; (c) $\dot x=-x+w/(1+x^2)$; (d) $\dot x=-x^3+x^3w$.
\end{problema}

\begin{problema}[Sontag robusto]\label{prob:av-sontagkappa}
\etiquetaM{} Para la planta $\dot x=x+(x^2+1)u$ con $V=\tfrac12x^2$: (a) escriba la ley de Sontag con $\kappa=1$; (b) muestre que el lazo es ISS respecto a un error aditivo $w$ en el actuador; (c) con \cmd{crearModelo\_RedisenoISS.m} y $w\equiv1$, compare con la cancelación exacta y con $k-\Lie gV$.
\end{problema}

\begin{problema}[ganancia de una cascada]\label{prob:av-gananciacascada}
Para $\dot x_1=-x_1^3+x_2$, $\dot x_2=-x_2+w$: (a) dé las ganancias ISS de cada subsistema; (b) con $\gamma_1(p+q)\le\gamma_1(2p)+\gamma_1(2q)$, dé una ganancia de la cascada respecto a $\abs{x_1}+\abs{x_2}$.
\end{problema}

\begin{problema}[una pequeña ganancia no lineal]\label{prob:av-pgnolineal}
\etiquetaEstrella{} Para $\dot x_1=-x_1+x_2^{1/3}$, $\dot x_2=-x_2+cx_1^3$ con $c>0$: (a) dé $\gamma_1$ y $\gamma_2$; (b) ¿para qué $c$ se cumple $\gamma_1\circ\gamma_2(r)<r$ para todo $r>0$? (c) Encuentre los equilibrios y explique por qué la respuesta de (b) no puede mejorarse.
\end{problema}

%% ---------------------------------------------------------------------
\section{Soluciones seleccionadas}
%% ---------------------------------------------------------------------

\paragraph{Problema~\ref{prob:av-cubica}}
(a) $\Vdot\le-x^4+\delta\abs x<0$ si $\abs x>\delta^{1/3}$; con $\theta$, $\mu=(\delta/\theta)^{1/3}$, y como $\alpha_1=\alpha_2=\tfrac12r^2$, $b=\mu$: cualquier $b>\delta^{1/3}$ es cota última. (b) Porque $V$ es ``redonda'': sus conjuntos de nivel son bolas. En el oscilador, $\lmax/\lmin\approx4$. (c) Con $\delta=0.5$, $\delta^{1/3}=0.794$; la simulación da $\limsup\abs x=0.49$: la senoidal no es la peor perturbación de amplitud $0.5$.

\paragraph{Problema~\ref{prob:av-esperaresorte}}
(a) $c=\tfrac32\cdot4+\tfrac12\cdot16=14$; $\mu=0.745$, $\eta=1.809\mu^2+\tfrac12\mu^4=1.16$; $k=(1-\theta)\mu^2=0.056$. (b) $T\le(c-\eta)/k=231$. (c) La simulación entra a $\Nivel\eta$ en $t\approx2.1$ con $w=0.3\cos t$ (entre $1.7$ y $2.5$ con $w\equiv\pm0.3$): la cota usa el ritmo del borde interior de la corona durante todo el recorrido (\cref{tr:av-T}).

\paragraph{Problema~\ref{prob:av-integrador}}
(a) $x(t)=x(t_0)+\delta(\cos t_0-\cos t)$, así que $\abs{x(t)}\le\abs{x(t_0)}+2\delta$. (b) Quien empieza en $x(t_0)=100$ se queda cerca de $100$. (c) No hay $W_3$ con $\Vdot\le-W_3$ fuera de una bola. Ninguna $V$ lo arregla: si existiera, el teorema daría UUB, que es falso. El sistema no tiene disipación.

\paragraph{Problema~\ref{prob:av-margen}}
(a) $P=\bigl[\begin{smallmatrix}1.25&0.25\\0.25&0.25\end{smallmatrix}\bigr]$. (b) $\lmax=1.309$, $c_4=2.618$ y $\nu^*=0.382$. (c) Con $g_p=\nu[0;\,x_1]$, la matriz $\bigl[\begin{smallmatrix}0&1\\-2+\nu&-3\end{smallmatrix}\bigr]$ deja de ser Hurwitz en $\nu=2$. Es mucho mayor que $\nu^*$: la garantía cubre todas las perturbaciones de ese tamaño, incluidas las variantes en el tiempo y no lineales.

\paragraph{Problema~\ref{prob:av-dobleintegrador}}
(a) Es $V=x\T Px$ con $q=\tfrac12>0$: por la \cref{tr:av-noclf}, es CLF. Explícitamente, $\ell_g=x_1+2x_2$ y en $x_2=-x_1/2$, $\ell_f=(2x_1+x_2)x_2=-\tfrac34x_1^2<0$. (b) $k=-(\ell_f+\sqrt{\ell_f^2+\ell_g^4})/\ell_g$ con $\ell_f=(2x_1+x_2)x_2$ cuadrático y $\ell_g$ lineal: $k(sx)/s\to-(\ell_f+\abs{\ell_f})/\ell_g$, que vale $-2\ell_f/\ell_g$ donde la deriva estorba y $0$ donde ayuda. (c) La pregunta es la trampa: el límite es continuo y homogéneo de grado uno, pero no lineal en $x$, así que la ley es continua (hay control pequeño) y \emph{no} es diferenciable en el origen. No hay polos que calcular. Una ley de Sontag puede ser continua sin ser $\Cont1$, aunque $V$, $f$ y $g$ sean polinomios.

\paragraph{Problema~\ref{prob:av-conservar}}
(a) $k=-x\bigl((1-x^2)+\sqrt{(1-x^2)^2+1}\bigr)$. (b) Sontag arranca con $u(0)=-0.19$ y gasta $\int u^2=1.46$; la ley que cancela arranca con $u(0)=21$ y gasta $58.5$. (c) Sontag también es más rápido: el lazo es $\dot x=-x\sqrt{(1-x^2)^2+1}$, con ritmo $\ge1$ en todo punto ($\sqrt2$ cerca del origen), y llega a $\abs x<0.01$ en $t=4.0$ contra $5.7$. No hay paradoja: la ley que cancela paga para quitar el término $-x^3$, que ayudaba, y luego impone un ritmo $1$; Sontag deja que $-x^3$ trabaje lejos del origen y solo empuja donde la deriva estorba.

\paragraph{Problema~\ref{prob:av-mrac}}
(a) $\dot e=a^*y-ky+r+ay_m-r=-ae-\tilde ky$. (b) $\Vdot=-ae^2-\tilde key+\tilde key=-ae^2\le0$: $e$ y $\tilde k$ acotadas; $y_m$ acotada porque $a>0$ y $r$ acotada; $y=e+y_m$ acotada; $\dot e$ acotada; $\ddot V=-2ae\dot e$ acotada; por Barbalat, $e\to0$. (c) No necesariamente: solo si $y$ es persistentemente excitante. Con $r\equiv0$ e $y\to0$, $k$ se congela en cualquier valor.

\paragraph{Problema~\ref{prob:av-sint2}}
(a) $\abs h\le1/t$. (b) $\dot h=2\cos(t^2)-\sin(t^2)/t^2$ oscila con amplitud cercana a $2$. (c) La continuidad uniforme de $\dot h$: su derivada contiene $-4t\sin(t^2)$, que no está acotada.

\paragraph{Problema~\ref{prob:av-issono}}
(a) No ($r=1<s=2$): con $w\equiv c$ y $x(0)>1/c$, $\dot x=x(cx-1)>0$ y escapa. (b) Sí ($r=3>s=2$): $\abs x\ge\abs w/\theta$ da $\Vdot\le-(1-\theta)x^4$, con $\gamma(r)=r/\theta$. (c) Sí: $\Vdot\le-x^2+\abs x\abs w$ y $\gamma(r)=r/\theta$; el factor $1/(1+x^2)\le1$ hace que la entrada entre a lo sumo linealmente. (d) No ($r=s=3$): con $w\equiv2$, $\dot x=x^3$ escapa.

\paragraph{Problema~\ref{prob:av-sontagkappa}}
(a) $\ell_f=x^2$, $\ell_g=x(x^2+1)$ y $k=-\bigl(1+\frac{x^2+\sqrt{x^4+x^4(x^2+1)^4}}{x^2(x^2+1)^2}\bigr)x(x^2+1)$. (b) $\Vdot\le-W-\tfrac12\ell_g^2+\tfrac12w^2$ con $W=x^2\sqrt{1+(x^2+1)^4}\ge\sqrt2x^2$, que es de clase $\Kinf$: ISS. (c) La cancelación escapa en $2.42$; las otras dos quedan acotadas ($k-\Lie gV$ se asienta en $0.4887$ y Sontag con $\kappa=1$ en $0.386$).

\paragraph{Problema~\ref{prob:av-gananciacascada}}
(a) $\gamma_1(r)=(r/\theta)^{1/3}$ (como $\dot x=-x^3+w$) y $\gamma_2(r)=r$. (b) $\abs{x_1}+\abs{x_2}$ queda, a la larga, por debajo de $\gamma_1(\gamma_2(r))+\gamma_2(r)$, y con la desigualdad del enunciado aplicada a los transitorios, $\gamma(r)=(2r/\theta)^{1/3}+r$.

\paragraph{Problema~\ref{prob:av-pgnolineal}}
(a) $\gamma_1(r)=r^{1/3}$ y $\gamma_2(r)=cr^3$ (los dos subsistemas son lineales estables con entradas $x_2^{1/3}$ y $cx_1^3$). (b) $\gamma_1(\gamma_2(r))=c^{1/3}r<r$ si y solo si $c<1$. (c) Los equilibrios cumplen $x_1=(cx_1^3)^{1/3}=c^{1/3}x_1$: con $c\neq1$ solo el origen; con $c=1$, todo punto de la curva $x_2=x_1^3$ es equilibrio, y el origen no es AS. La condición no se puede mejorar en $c=1$.

\respuestaspausas

